% =============================================================================
% Emergent Constitutions: A Heterogeneous-Agent Model with Endogenous
% Governance and LLM-Augmented Political Decision Making
% =============================================================================
\documentclass[12pt]{article}

% --- Packages ---
\usepackage[utf8]{inputenc}
\usepackage[T1]{fontenc}
\usepackage{amsmath,amssymb,amsthm}
\usepackage{mathtools}
\usepackage[numbers,sort&compress]{natbib}
\usepackage{booktabs}
\usepackage[ruled,vlined,linesnumbered]{algorithm2e}
\usepackage{geometry}
\usepackage{hyperref}
\usepackage{cleveref}
\usepackage{enumitem}
\usepackage{graphicx}
\usepackage{setspace}
\usepackage{caption}
\usepackage{subcaption}
\usepackage{xcolor}

% --- Page geometry ---
\geometry{margin=1in}
\onehalfspacing

% --- Theorem environments ---
\newtheorem{definition}{Definition}
\newtheorem{proposition}{Proposition}
\newtheorem{assumption}{Assumption}
\newtheorem{remark}{Remark}

% --- Custom commands ---
\newcommand{\E}{\mathbb{E}}
\newcommand{\R}{\mathbb{R}}
\newcommand{\N}{\mathbb{N}}
\newcommand{\calC}{\mathcal{C}}
\newcommand{\calR}{\mathcal{R}}
\newcommand{\calS}{\mathcal{S}}
\newcommand{\calA}{\mathcal{A}}
\newcommand{\calF}{\mathcal{F}}
\newcommand{\calM}{\mathcal{M}}

% --- Metadata ---
\title{Emergent Constitutions: A Heterogeneous-Agent Model with Endogenous Governance and LLM-Augmented Political Decision Making}

\author{%
  Antonio Mele \\
  \textit{Working Paper}
}

\date{\today}

% =============================================================================
\begin{document}
% =============================================================================

\maketitle

% --- Abstract ---
\begin{abstract}
We develop a Heterogeneous Agent New Keynesian (HANK) model with endogenous governance institutions. Agents choose between working and entrepreneurship via Bellman-based occupational choice, face idiosyncratic productivity and ability shocks discretized via the Rouwenhorst method, hold two-asset portfolios of liquid bonds and illiquid capital with convex adjustment costs, and collectively design taxation, redistribution, public goods provision, regulation, and voting rules through a mutable constitutional framework. Household optimization is solved via the Endogenous Grid Method (EGM) of Carroll (2006), with Fella (2014) upper envelope for non-monotonicities. Walrasian market clearing with heterogeneous firms uses bisection on excess labor demand. The cross-sectional distribution is tracked via the Kolmogorov Forward Equation using the Young (2010) lottery allocation method. A government sector issues bonds, finances spending through a budget constraint, and applies fiscal rules for debt sustainability. Nominal rigidities follow the Rotemberg (1982) sticky price framework with a Taylor rule and Fisher equation. Political preferences are microfounded: a political utility function enters the Bellman equation as a constant flow bonus, yielding structurally consistent proposal evaluation and voting. Political deliberation is delegated to large language models (LLMs) that receive Bellman-derived context, mimicking bounded rationality in institutional design. We characterize the recursive competitive equilibrium, describe the computational algorithm, and discuss implications for the co-evolution of institutions and inequality.
\end{abstract}

\medskip

\noindent\textbf{Keywords:} HANK, heterogeneous agents, endogenous grid method, New Keynesian, endogenous institutions, occupational choice, LLM agents, constitutional dynamics, two-asset model, Kolmogorov Forward Equation

\medskip

\noindent\textbf{JEL Classification:} D72, E21, E31, E52, H11, H63, L26, C63

\newpage
\tableofcontents
\newpage

% =============================================================================
\section{Introduction}\label{sec:intro}
% =============================================================================

A central question in political economy is how institutions emerge, persist, and evolve. Standard macroeconomic models take institutions---tax codes, property rights, voting rules---as exogenous parameters. Yet real-world governance is the product of collective action by heterogeneous agents whose economic circumstances shape their political preferences and vice versa. This paper develops a computational framework that endogenizes the entire governance layer within a heterogeneous-agent macroeconomic model.

\subsection{Motivation}

The canonical incomplete-markets models of \citet{aiyagari1994}, \citet{huggett1993}, and \citet{krusellsmith1998} have yielded deep insights into wealth inequality, precautionary savings, and aggregate dynamics under idiosyncratic risk. However, these models typically hold fiscal policy fixed: the tax rate, transfer schedule, and public goods provision are parameters, not outcomes. In reality, agents vote on these policies, propose new rules, and form coalitions to advance their interests. The feedback loop---from economic inequality to political action to institutional change to new equilibrium outcomes---is precisely what this paper formalizes.

\subsection{Contributions}

This paper makes five contributions.

\textbf{First}, we build a DSGE-HA model with occupational choice between workers and entrepreneurs, endogenous firm dynamics with R\&D-driven productivity growth, and Cobb-Douglas production at the firm level. Idiosyncratic productivity and entrepreneurial ability follow AR(1) processes discretized via the Rouwenhorst method \citep{kopeckysuen2010}, while aggregate TFP follows a separate AR(1) process. Household decisions are solved via the Endogenous Grid Method (EGM) of \citet{carroll2006} on the Bellman equation, with an intratemporal first-order condition for leisure and Fella (2014) upper envelope for non-monotonicities.

\textbf{Second}, we introduce a mutable constitutional framework with eight rule types---tax schedules, transfer programs, public goods provision, market regulation, firm regulation, voting procedures, property rights, and custom rules---each with sandboxed enforcement code. Constitutional rules feed back into agent budget constraints through taxation, transfers, and regulation. Agents propose, validate, vote on, and amend rules through a formal political process in which voting thresholds themselves are endogenous.

\textbf{Third}, we extend the model to a full HANK framework. Households hold two-asset portfolios of liquid government bonds and illiquid capital with convex adjustment costs \citep{kaplanmollviolante2018}. A government sector issues bonds, finances spending through a budget constraint, and applies fiscal rules for debt sustainability. Nominal rigidities follow the \citet{rotemberg1982} sticky price framework, with a Taylor rule for monetary policy and the Fisher equation linking nominal and real rates. Market clearing uses Walrasian bisection on excess labor demand with heterogeneous firms, and the cross-sectional distribution is tracked via the Kolmogorov Forward Equation using the \citet{young2010} lottery allocation method.

\textbf{Fourth}, we microfound political preferences within the Bellman equation. A political utility function $v(\calC; \theta)$---depending on the constitution and the agent's ideological value vector---enters the household's value function as a constant flow bonus, yielding structurally consistent proposal evaluation and voting decisions derived from the value function.

\textbf{Fifth}, we integrate large language models (LLMs) into the political decision layer. Economic optimization is analytically well-defined (the Bellman equation has a known structure), so we solve it numerically. Political deliberation, by contrast, is open-ended: what tax rate to propose, how to frame a regulation, whether a means-tested transfer is better than a universal one. We delegate these decisions to LLMs that receive Bellman-derived context about the economy and society, producing proposals and votes that reflect bounded rationality in institutional design. A benchmark mode with purely numerical governance enables controlled comparisons.

\subsection{Related Literature}

Our model builds on several strands of the literature. The heterogeneous-agent tradition originates with \citet{bewley1977}, \citet{huggett1993}, and \citet{aiyagari1994}, with \citet{krusellsmith1998} extending to aggregate shocks. \citet{cagettideNardi2006} introduce entrepreneurship into the Bewley framework. We follow their approach of modeling occupational choice with heterogeneous entrepreneurial ability but add endogenous governance and LLM-driven institutional dynamics.

On the computational side, our discretization of AR(1) processes follows \citet{kopeckysuen2010}, who show that the Rouwenhorst method dominates Tauchen's method for highly persistent processes. Our household solver uses the Endogenous Grid Method (EGM) of \citet{carroll2006}, which avoids the inner maximization loop of VFI by inverting the Euler equation on an exogenous savings grid and recovering the endogenous current asset level. We apply the Fella (2014) upper envelope to handle non-monotonicities near kinks in the budget constraint.

The HANK literature originates with \citet{kaplanmollviolante2018}, who show that the interaction of liquid and illiquid assets is crucial for understanding the transmission of monetary policy. Our two-asset extension follows their framework with convex adjustment costs for illiquid capital deposits. The Kolmogorov Forward Equation for distribution tracking uses the \citet{young2010} lottery allocation method for mass conservation on a discrete grid. The nominal block follows \citet{rotemberg1982} for sticky prices, with a Taylor rule satisfying the Taylor principle.

The endogenous institutions literature includes \citet{acemoglurobinson2006}, who study the dynamics of political institutions and economic outcomes, and \citet{acemoglu2005}, who formalize the relationship between political power and economic institutions. Our model operationalizes these ideas computationally: agents with heterogeneous preferences over equality and liberty vote on the very rules that govern economic outcomes.

The use of LLMs as economic agents is nascent. We contribute by showing how LLM-generated political decisions can be integrated into a fully specified general equilibrium model while preserving determinism (temperature zero, seeded RNG) and providing analytical fallback (benchmark mode). A novel element is our microfounded political utility, which enters the Bellman equation and provides structurally consistent preferences that inform LLM-generated political decisions.

\medskip

The remainder of the paper is organized as follows. \Cref{sec:environment} describes the model environment. \Cref{sec:governance} formalizes the constitutional framework. \Cref{sec:equilibrium} defines the recursive competitive equilibrium. \Cref{sec:computation} details the computational method, including the EGM solver and Bellman-based occupational choice. \Cref{sec:distribution} presents the KFE distribution tracking. \Cref{sec:two_asset} describes the two-asset household structure. \Cref{sec:government} develops the government sector and bond market. \Cref{sec:nominal} introduces nominal rigidities. \Cref{sec:political_utility} develops the microfounded political utility framework. \Cref{sec:llm} discusses the LLM integration. \Cref{sec:calibration} presents the calibration, including the SMM procedure. \Cref{sec:discussion} discusses results and implications. \Cref{sec:conclusion} concludes.


% =============================================================================
\section{Model Environment}\label{sec:environment}
% =============================================================================

\subsection{Time, Agents, and State Space}\label{sec:state}

Time is discrete, indexed by $t = 0, 1, 2, \ldots, T$. The economy is populated by $N \geq 20$ agents, indexed by $i = 1, \ldots, N$.

\begin{definition}[Individual State]
Each agent $i$ at time $t$ is characterized by the individual state vector
\begin{equation}\label{eq:individual_state}
  s_{it} = (a_{it},\, z_{it},\, e_{it},\, o_{it}),
\end{equation}
where $a_{it} \in [\underline{a}, \infty)$ denotes assets (wealth), $z_{it} \in \calS_z$ is idiosyncratic labor productivity, $e_{it} \in \calS_e$ is entrepreneurial ability, and $o_{it} \in \{W, E\}$ denotes occupation (worker or entrepreneur).
\end{definition}

\begin{definition}[Aggregate State]
The aggregate state at time $t$ is
\begin{equation}\label{eq:aggregate_state}
  \Gamma_t = (\mu_t,\, \calC_t,\, w_t,\, r_t,\, A_t),
\end{equation}
where $\mu_t$ is the cross-sectional distribution of individual states, $\calC_t$ is the constitution (a mutable set of rules), $w_t$ is the wage, $r_t$ is the interest rate, and $A_t$ is aggregate total factor productivity (TFP).
\end{definition}


\subsection{Preferences}\label{sec:preferences}

Agents have Cobb-Douglas preferences over private consumption $c$, leisure $\ell$, and per-capita public goods $G$:

\begin{equation}\label{eq:utility}
  u(c, \ell, G) = c^{\alpha_u} \cdot \ell^{\beta_u} \cdot G^{\gamma_u}, \qquad \alpha_u + \beta_u + \gamma_u = 1,
\end{equation}
where $\alpha_u, \beta_u, \gamma_u > 0$ are the weights on consumption, leisure, and public goods, respectively. The constraint $\alpha_u + \beta_u + \gamma_u = 1$ ensures homogeneity of degree one.

The lifetime objective is
\begin{equation}\label{eq:lifetime}
  \max \; \E_0 \sum_{t=0}^{\infty} \beta_d^{\, t} \, u(c_{it}, \ell_{it}, G_t),
\end{equation}
where $\beta_d \in (0, 1)$ is the subjective discount factor.

\begin{remark}[Preference Heterogeneity]
The model supports two modes. Under \emph{homogeneous preferences}, all agents share the same $(\alpha_u, \beta_u, \gamma_u, \beta_d)$, enabling a significant computational speedup: VFI is solved once and policies are interpolated per agent. Under \emph{heterogeneous preferences}, each agent draws utility weights from a Dirichlet distribution and $\beta_d$ from $U[0.9, 0.99]$, requiring per-agent VFI.
\end{remark}

In addition to economic preferences, each agent possesses an ideological value vector $v_i = (v_i^{eq}, v_i^{lib})$ with $v_i^{eq} + v_i^{lib} = 1$, where $v_i^{eq}$ represents the preference weight toward equality and $v_i^{lib}$ toward liberty. These values influence political decisions---agents with high $v_i^{eq}$ tend to favor redistribution, while those with high $v_i^{lib}$ favor low taxes and free markets.


\subsection{Technology}\label{sec:technology}

\subsubsection{Workers}

A worker with productivity $z_{it}$ and leisure choice $\ell_{it} \in [0, 1]$ supplies effective labor
\begin{equation}\label{eq:effective_labor}
  h_{it} = z_{it} \cdot (1 - \ell_{it}),
\end{equation}
and earns labor income $w_t \cdot h_{it}$.

\subsubsection{Entrepreneurs and Firms}

An entrepreneur $i$ owns a firm $f$ that produces output according to the Cobb-Douglas production function
\begin{equation}\label{eq:firm_production}
  Y_f = A_f \cdot K_f^{\alpha} \cdot L_f^{1-\alpha},
\end{equation}
where $A_f = e_{it} \cdot A_t$ is firm-specific TFP (the product of the entrepreneur's ability and aggregate TFP), $K_f$ is the firm's capital, $L_f$ is effective labor hired, and $\alpha \in (0, 1)$ is the capital share.

The first-order condition for optimal labor demand is
\begin{equation}\label{eq:foc_labor}
  (1 - \alpha) \cdot A_f \cdot K_f^{\alpha} \cdot L_f^{-\alpha} = w_t,
\end{equation}
yielding optimal labor
\begin{equation}\label{eq:optimal_labor}
  L_f^* = \left(\frac{(1-\alpha) \cdot A_f \cdot K_f^{\alpha}}{w_t}\right)^{1/\alpha}.
\end{equation}

Firm profit is
\begin{equation}\label{eq:firm_profit}
  \pi_f = Y_f - w_t L_f^* - (r_t + \delta) K_f,
\end{equation}
where $\delta$ is the depreciation rate and $r_t + \delta$ is the rental rate of capital.

\subsubsection{R\&D and Firm-Specific TFP Growth}

Active firms may invest in R\&D at expenditure $x_f^{rd}$. The probability of a TFP improvement is
\begin{equation}\label{eq:rd_prob}
  p_f = p_0 \cdot \sqrt{\frac{x_f^{rd}}{\bar{x}^{rd}}},
\end{equation}
where $p_0$ is the base success probability and $\bar{x}^{rd}$ is the mean R\&D spending across all active firms. Upon success, firm TFP updates as
\begin{equation}\label{eq:rd_improvement}
  A_f' = A_f \cdot (1 + \eta), \qquad \eta \sim \max\left\{0,\, \mathcal{N}(\mu_{rd}, \sigma_{rd}^2)\right\}.
\end{equation}


\subsection{Occupational Choice: Entry and Exit}\label{sec:entry_exit}

Occupational choice is determined by comparing lifetime utility as a worker versus an entrepreneur. This utility-based framework replaces profit-based entry/exit conditions, capturing the full cost of entrepreneurship including reduced leisure.

\subsubsection{Firm Value}

The value of a firm to an entrepreneur with ability $e$ and capital $K$ uses a closed-form infinite-horizon formula based on the stationary distribution of the ability Markov chain:
\begin{equation}\label{eq:firm_value}
  V(e, K) = \pi(e, K) + \frac{\beta_f \cdot \pi(\bar{e}, K)}{1 - \beta_f},
\end{equation}
where $\pi(e, K)$ is current-period profit at actual ability $e$, $\pi(\bar{e}, K)$ is the stationary expected profit evaluated at $\bar{e} = \E_{\pi^*}[e]$ (the mean ability under the stationary distribution $\pi^*$ of the ability Markov chain), and $\beta_f$ is the firm discount factor. The first term reflects exact current-period profit; the second term treats future profits as a perpetuity at the stationary expected level.

The stationary distribution $\pi^*$ is computed once at initialization via power iteration on the ability transition matrix $\Pi_e$, converging to $\|\pi^{(k+1)} - \pi^{(k)}\|_\infty < 10^{-12}$.

\subsubsection{Worker Lifetime Value}

The lifetime utility value of remaining a worker is approximated as a perpetuity of per-period utility at the household's optimal allocation:
\begin{equation}\label{eq:worker_value}
  V^W(s_i) = \frac{u(c^*_W, \ell^*_W, G)}{1 - \beta_d},
\end{equation}
where $\ell^*_W \approx 0.35$ is the typical optimal leisure from VFI, and $c^*_W \approx 0.7 \times y_{avail}$ is approximate optimal consumption (consuming roughly 70\% of available income including labor and asset income).

\subsubsection{Entrepreneur Lifetime Value}

The lifetime utility value of entrepreneurship accounts for firm profits, reduced leisure from running a firm, and the entry cost penalty:
\begin{equation}\label{eq:entre_value}
  V^E(s_i) = \frac{u(c^*_E, \ell_E, G)}{1 - \beta_d} - \Phi(F),
\end{equation}
where $\ell_E = 0.2$ reflects the reduced leisure of an entrepreneur (running a firm takes effort), $c^*_E$ is consumption from firm profits plus asset income, and $\Phi(F)$ is the entry cost expressed in utility units:
\begin{equation}\label{eq:entry_cost_utility}
  \Phi(F) = F \cdot \frac{\alpha_u}{c^*_E} \cdot u(c^*_E, \ell_E, G),
\end{equation}
i.e., the entry cost $F$ scaled by the marginal utility of consumption. The term $\Phi(F)$ is included only for new entrants, not for continuing entrepreneurs.

\subsubsection{Entry Condition}

A worker enters entrepreneurship if two conditions hold:
\begin{equation}\label{eq:entry}
  V^E(s_i) > V^W(s_i) \qquad \text{and} \qquad a_{it} \geq \underline{K} + F,
\end{equation}
where $\underline{K}$ is the minimum capital requirement and $F$ is the sunk entry cost. The first condition ensures that entrepreneurship yields higher lifetime utility; the second ensures financial feasibility.

\subsubsection{Exit Condition}

An existing entrepreneur exits (reverts to worker) if
\begin{equation}\label{eq:exit}
  V^W(s_i) > V^E(s_i),
\end{equation}
where $V^E$ is evaluated without the entry cost penalty (since the entrepreneur is already operating). Exit occurs when lifetime worker utility exceeds the ongoing value of the firm.

\subsubsection{Optimal Capital}

The optimal capital level $K^*$ maximizes firm value subject to feasibility:
\begin{equation}\label{eq:optimal_capital}
  K^* = \arg\max_{K \in [\underline{K}, \, a_{it} - F]} V(e_i, K),
\end{equation}
solved via grid search over 20 candidate capital levels.


\subsection{Stochastic Processes}\label{sec:shocks}

The model features three sources of uncertainty.

\subsubsection{Idiosyncratic Productivity}

Labor productivity follows an AR(1) process in logs:
\begin{equation}\label{eq:ar1_z}
  \log z_{it} = \rho_z \log z_{i,t-1} + \varepsilon_{it}^z, \qquad \varepsilon_{it}^z \sim \mathcal{N}(0, \sigma_z^2).
\end{equation}

\subsubsection{Entrepreneurial Ability}

Entrepreneurial ability follows an independent AR(1) process:
\begin{equation}\label{eq:ar1_e}
  \log e_{it} = \rho_e \log e_{i,t-1} + \varepsilon_{it}^e, \qquad \varepsilon_{it}^e \sim \mathcal{N}(0, \sigma_e^2).
\end{equation}

\subsubsection{Aggregate TFP}

Aggregate TFP follows an AR(1) process:
\begin{equation}\label{eq:ar1_A}
  \log A_t = \rho_A \log A_{t-1} + \varepsilon_t^A, \qquad \varepsilon_t^A \sim \mathcal{N}(0, \sigma_A^2).
\end{equation}

\subsubsection{Rouwenhorst Discretization}

The continuous AR(1) processes \eqref{eq:ar1_z} and \eqref{eq:ar1_e} are discretized into finite-state Markov chains using the Rouwenhorst method \citep{kopeckysuen2010}. For a process with persistence $\rho$ and innovation volatility $\sigma$:

\begin{enumerate}[label=\textit{Step \arabic*.}]
  \item Compute the unconditional standard deviation: $\sigma_z^{unc} = \sigma / \sqrt{1 - \rho^2}$.
  \item Set grid bounds: $z_{max} = \sigma_z^{unc} \cdot \sqrt{n-1}$, where $n$ is the number of grid points.
  \item Construct an equally-spaced log grid: $\{-z_{max}, -z_{max} + \Delta, \ldots, z_{max}\}$ with $\Delta = 2 z_{max}/(n-1)$.
  \item Set the transition probability parameter: $p = (1 + \rho)/2$.
  \item Build the transition matrix recursively. The base case ($n = 2$) is
  \begin{equation}\label{eq:rouwenhorst_base}
    \Pi_2 = \begin{pmatrix} p & 1-p \\ 1-p & p \end{pmatrix}.
  \end{equation}
  For $n > 2$, construct $\Pi_n$ from $\Pi_{n-1}$ via four-quadrant expansion:
  \begin{equation}\label{eq:rouwenhorst_recursion}
    \widetilde{\Pi}_n = p \begin{pmatrix} \Pi_{n-1} & \mathbf{0} \\ \mathbf{0}^\top & 0 \end{pmatrix}
    + (1-p) \begin{pmatrix} \mathbf{0} & \Pi_{n-1} \\ 0 & \mathbf{0}^\top \end{pmatrix}
    + (1-p) \begin{pmatrix} \mathbf{0}^\top & 0 \\ \Pi_{n-1} & \mathbf{0} \end{pmatrix}
    + p \begin{pmatrix} 0 & \mathbf{0}^\top \\ \mathbf{0} & \Pi_{n-1} \end{pmatrix},
  \end{equation}
  followed by row normalization: $\Pi_n(i, \cdot) = \widetilde{\Pi}_n(i, \cdot) / \sum_j \widetilde{\Pi}_n(i, j)$.
  \item Convert the log grid to levels: $\calS = \{\exp(z_1), \ldots, \exp(z_n)\}$.
\end{enumerate}

The stationary distribution $\pi^*$ of $\Pi$ is computed by power iteration until $\|\pi^{(k+1)} - \pi^{(k)}\|_\infty < 10^{-12}$.


\subsection{Markets and Budget Constraints}\label{sec:markets}

\subsubsection{Factor Market Clearing}

In equilibrium, factor markets clear. Labor market clearing requires
\begin{equation}\label{eq:labor_clearing}
  \sum_{f \in \calF_t} L_f^*(w_t, r_t) = \sum_{i: o_{it} = W} z_{it} \cdot (1 - \ell_{it}),
\end{equation}
where the left-hand side is aggregate labor demand from firms and the right-hand side is aggregate effective labor supply from workers.

Capital market clearing requires
\begin{equation}\label{eq:capital_clearing}
  \sum_{f \in \calF_t} K_f = \sum_{i=1}^{N} a_{it},
\end{equation}
where the left-hand side is aggregate capital demand (firm capital) and the right-hand side is aggregate savings.

\subsubsection{Worker Budget Constraint}

A worker faces the budget constraint
\begin{equation}\label{eq:worker_budget}
  a_{i,t+1} = (1 + r_t) a_{it} + w_t z_{it} (1 - \ell_{it}) - c_{it} - T(y_{it}) + Tr_{it},
\end{equation}
where $T(y_{it})$ is the tax schedule evaluated at labor income $y_{it} = w_t z_{it}(1 - \ell_{it})$, and $Tr_{it}$ is transfers received. Agents face the borrowing constraint $a_{i,t+1} \geq \underline{a}$.

\subsubsection{Entrepreneur Budget Constraint}

An entrepreneur additionally receives firm profit:
\begin{equation}\label{eq:entre_budget}
  a_{i,t+1} = (1 + r_t) a_{it} + w_t z_{it}(1 - \ell_{it}) + \pi_{f(i)} - c_{it} - T(y_{it}) + Tr_{it},
\end{equation}
where $\pi_{f(i)}$ is the profit of the firm owned by agent $i$.

\subsubsection{Goods Market Identity}

By Walras' law, the goods market clears:
\begin{equation}\label{eq:goods_clearing}
  Y_t = C_t + I_t + G_t,
\end{equation}
where $Y_t = \sum_f Y_f$ is aggregate output, $C_t = \sum_i c_{it}$ is aggregate consumption, $I_t = \sum_f \delta K_f$ is replacement investment, and $G_t$ is government (public goods) spending.


% =============================================================================
\section{Governance and Constitutional Framework}\label{sec:governance}
% =============================================================================

\subsection{Constitutional State Space}\label{sec:constitution_state}

The constitution at time $t$ is a finite collection of rules:
\begin{equation}\label{eq:constitution}
  \calC_t = \{R_1, R_2, \ldots, R_M\},
\end{equation}
where each rule $R_m = (name_m, type_m, \theta_m, f_m(\cdot))$ consists of a unique identifier, a type from the set
\begin{equation}\label{eq:rule_types}
  \mathcal{T} = \{\text{tax\_schedule},\, \text{transfer\_program},\, \text{public\_goods},\, \text{market\_regulation},\, \text{firm\_regulation},\, \text{voting\_procedure},\, \text{property\_rights},\, \text{custom}\},
\end{equation}
a parameter vector $\theta_m$ (e.g., $\theta = \{rate: 0.2\}$ for a tax rule), and an enforcement function $f_m(\cdot)$ specified as a sandboxed expression.


\subsection{Fiscal Policy}\label{sec:fiscal}

\subsubsection{Taxation}

The tax on agent $i$'s income is determined by all active rules of type \texttt{tax\_schedule}:
\begin{equation}\label{eq:taxation}
  T_i = \sum_{m: type_m = \text{tax\_schedule}} f_m(y_i; \theta_m),
\end{equation}
subject to the non-negativity constraint $T_i \in [0, y_i]$ (taxes cannot exceed income). The default rule is a flat tax: $T_i = \tau \cdot y_i$ with $\tau \in [0, 1]$.

Total tax revenue is $R = \sum_{i=1}^{N} T_i$.

\subsubsection{Public Goods}

Public goods spending is determined by rules of type \texttt{public\_goods}:
\begin{equation}\label{eq:public_goods}
  G_t = \frac{\phi \cdot R}{N},
\end{equation}
where $\phi \in [0, 1]$ is the fraction of revenue allocated to public goods, specified by the constitutional rule parameters.

\subsubsection{Transfers}

The remaining revenue $R - \phi R$ is distributed according to transfer rules. Three transfer methods are available.

Under \emph{equal-share} redistribution:
\begin{equation}\label{eq:equal_transfer}
  Tr_i = \frac{(1 - \phi) R}{N}.
\end{equation}
Under \emph{means-tested} redistribution:
\begin{equation}\label{eq:means_tested}
  Tr_i = \frac{(1 - \phi) R \cdot \omega_i}{\sum_{j=1}^{N} \omega_j}, \qquad \omega_i = \frac{1}{\max\{a_{it}, 1\}},
\end{equation}
so that poorer agents receive larger transfers.

Under \emph{progressive} redistribution, the same inverse-wealth weighting as means-tested is used:
\begin{equation}\label{eq:progressive_transfer}
  Tr_i^{prog} = \frac{(1 - \phi) R \cdot \omega_i}{\sum_{j=1}^{N} \omega_j}, \qquad \omega_i = \frac{1}{\max\{a_{it}, 1\}}.
\end{equation}
The progressive method serves as an alternative label for the inverse-wealth weighting scheme, available as a distinct policy choice in the constitutional framework.


\subsection{Political Process}\label{sec:political}

The constitutional amendment process follows four stages.

\begin{enumerate}[label=\textit{Stage \arabic*.}]
  \item \textbf{Proposal.} At proposal-interval periods (every $K_{prop}$ periods), each agent may propose one constitutional change. A proposal specifies an action $\in \{$add, modify, remove$\}$, a rule name, a rule type (for additions), parameter values, and optionally enforcement code.

  \item \textbf{Validation.} Each proposal is validated against structural constraints:
  \begin{itemize}
    \item Tax rates must lie in $[0, 1]$.
    \item Transfer methods must be recognized (equal-share, means-tested, progressive).
    \item Public goods fractions must lie in $[0, 1]$.
    \item Enforcement code must pass an AST-level safety check (restricted to arithmetic, comparisons, and safe builtins; no imports, function definitions, or attribute access).
    \item Additions cannot duplicate existing rule names; modifications require existing rules; the active voting rule cannot be removed.
  \end{itemize}
  Invalid proposals are rejected before reaching the voting stage.

  \item \textbf{Voting.} Each agent casts a vote (for or against) on each validated proposal. The outcome is determined by the currently active voting procedure rule with threshold $\bar{\theta} \in (0, 1]$:
  \begin{equation}\label{eq:voting}
    \text{Proposal passes} \iff
    \begin{cases}
      V_{for} > N \cdot \bar{\theta}, & \text{if } \bar{\theta} \leq 0.5 \text{ (majority)}, \\[4pt]
      V_{for} \geq \lceil N \cdot \bar{\theta} \rceil, & \text{if } 0.5 < \bar{\theta} < 1 \text{ (supermajority)}, \\[4pt]
      V_{for} = N, & \text{if } \bar{\theta} = 1 \text{ (unanimity)},
    \end{cases}
  \end{equation}
  where $V_{for}$ is the number of votes in favor and $N$ is the total number of eligible voters. Ties are resolved in favor of the status quo.

  \item \textbf{Amendment.} Passed proposals are applied to the constitution in order. Each application re-validates the resulting rule before activation. The constitution is immutable within a period and only updated between periods.
\end{enumerate}


\subsection{Endogenous Feedback}\label{sec:feedback}

The distinguishing feature of this framework is the feedback loop between governance and economics. Constitutional rules determine:
\begin{itemize}
  \item The tax function $T(\cdot)$ that enters budget constraints \eqref{eq:worker_budget}--\eqref{eq:entre_budget};
  \item The transfer schedule $Tr_i$ that enters household income;
  \item Public goods $G_t$ that enters the utility function \eqref{eq:utility};
  \item Market and firm regulations that constrain economic behavior.
\end{itemize}
These in turn affect prices $(w_t, r_t)$, wealth dynamics $a_{i,t+1}$, occupational choice, and ultimately the political preferences that generate new proposals. The constitution is thus endogenous in the deepest sense: it is both a determinant and a consequence of economic equilibrium.


% =============================================================================
\section{Equilibrium Definition}\label{sec:equilibrium}
% =============================================================================

\subsection{Worker's Problem}

A worker with state $(a, z)$, facing prices $(w, r)$, public goods $G$, tax function $T(\cdot)$, and transfer $Tr$, solves the Bellman equation
\begin{equation}\label{eq:bellman_worker}
  V^W(a, z) = \max_{c, \ell} \left\{ u(c, \ell, G) + \beta_d \, \E\!\left[V(a', z') \,\big|\, z\right] \right\}
\end{equation}
subject to
\begin{align}
  a' &= (1 + r) a + w z (1 - \ell) - c - T(wz(1-\ell)) + Tr, \label{eq:bellman_budget}\\
  a' &\geq \underline{a}, \label{eq:bellman_borrow}\\
  c &\geq 0, \qquad \ell \in [0, 1], \label{eq:bellman_bounds}
\end{align}
where $V(a', z') = \max\{V^W(a', z'), V^E(a', z', e')\}$ incorporates the occupational choice.

\subsection{Entrepreneur's Problem}

An entrepreneur with state $(a, z, e)$ and firm capital $K$ solves
\begin{equation}\label{eq:bellman_entre}
  V^E(a, z, e) = \max_{c, \ell, K} \left\{ u(c, \ell, G) + \beta_d \, \E\!\left[V(a', z', e')\right] \right\}
\end{equation}
subject to the entrepreneur budget constraint \eqref{eq:entre_budget} and the entry/exit conditions of \Cref{sec:entry_exit}.

\subsection{Market Clearing}

\begin{definition}[Recursive Competitive Equilibrium]
A \emph{recursive competitive equilibrium} (RCE) consists of:
\begin{enumerate}
  \item Value functions $V^W(a, z)$ and $V^E(a, z, e)$;
  \item Policy functions $c(a, z)$, $\ell(a, z)$, and occupational choice $o(a, z, e)$;
  \item Price functions $w(\Gamma)$ and $r(\Gamma)$;
  \item A law of motion for the cross-sectional distribution $\mu_{t+1} = \Phi(\mu_t, \calC_t)$;
  \item A law of motion for the constitution $\calC_{t+1} = \Psi(\calC_t, \mu_t, w_t, r_t)$;
\end{enumerate}
such that:
\begin{enumerate}[label=(\alph*)]
  \item Given prices and the constitution, agents optimize (solve their respective Bellman equations);
  \item Factor markets clear (\Cref{eq:labor_clearing,eq:capital_clearing});
  \item The goods market identity \eqref{eq:goods_clearing} holds;
  \item The distribution evolves consistently with individual decisions and shocks;
  \item The constitution evolves through the political process of \Cref{sec:political}, given the economic state.
\end{enumerate}
\end{definition}

The key departure from standard RCE definitions is condition (e): the constitution is a state variable with an endogenous law of motion determined by agent proposals and votes.


% =============================================================================
\section{Computational Method}\label{sec:computation}
% =============================================================================

\subsection{Value Function Iteration}\label{sec:vfi}

The worker's Bellman equation \eqref{eq:bellman_worker} is solved numerically via VFI on a discretized state space.

\textbf{Asset grid.} We construct an exponentially-spaced grid of $N_a = 100$ asset levels on $[\underline{a}, a_{max}]$ with $a_{max} = 1000$:
\begin{equation}\label{eq:asset_grid}
  a_j = \underline{a} + (a_{max} - \underline{a}) \cdot \left(\frac{j}{N_a - 1}\right)^2, \qquad j = 0, 1, \ldots, N_a - 1.
\end{equation}
The quadratic spacing concentrates grid points near the borrowing constraint $\underline{a}$, where the policy function has the highest curvature.

\textbf{Leisure grid.} A uniform grid of $N_\ell = 21$ leisure values: $\ell_k = k/20$ for $k = 0, 1, \ldots, 20$.

\textbf{Consumption optimization.} For each $(a_j, z_s, \ell_k)$ triplet, the budget constraint determines the maximum feasible consumption $c_{max} = budget - \underline{a}$ where $budget = (1+r)a_j + wz_s(1-\ell_k) - T(wz_s(1-\ell_k)) + Tr$. We search over 11 equally-spaced consumption levels $c \in \{0, c_{max}/10, \ldots, c_{max}\}$.

\textbf{Expected continuation value.} The conditional expectation $\E[V(a', z') \mid z_s]$ is computed as
\begin{equation}\label{eq:ev}
  \E[V(a', z') \mid z_s] = \sum_{s'=1}^{N_z} \Pi(s, s') \cdot V^{interp}(a', s'),
\end{equation}
where $\Pi$ is the Rouwenhorst transition matrix and $V^{interp}$ is the value function linearly interpolated on the asset grid.

\textbf{Convergence.} VFI iterates until $\|V_{n+1} - V_n\|_\infty < \epsilon_{VFI}$ with $\epsilon_{VFI} = 10^{-6}$, or a maximum of 500 iterations.

\begin{algorithm}[t]
\caption{Value Function Iteration}\label{alg:vfi}
\KwIn{Prices $(w, r)$, public goods $G$, tax function $T(\cdot)$, transfer $Tr$, transition matrix $\Pi$}
\KwOut{Policy functions $c^*(a, z)$, $\ell^*(a, z)$}
Initialize $V^0(a_j, z_s)$ for all $(j, s)$\;
\For{$n = 1, \ldots, 500$}{
  \For{$j = 1, \ldots, N_a$; $\; s = 1, \ldots, N_z$}{
    $V^{best} \leftarrow -\infty$\;
    \For{$k = 0, \ldots, N_\ell$}{
      $\ell \leftarrow k / N_\ell$\;
      $budget \leftarrow (1+r) a_j + w z_s (1-\ell) - T(w z_s (1-\ell)) + Tr$\;
      $c_{max} \leftarrow \max(0,\; budget - \underline{a})$\;
      \For{$m = 0, \ldots, 10$}{
        $c \leftarrow c_{max} \cdot m/10$\;
        $a' \leftarrow \max(budget - c,\; \underline{a})$\;
        $EV \leftarrow \sum_{s'} \Pi(s,s') \cdot V^{n-1}_{interp}(a', s')$\;
        $val \leftarrow u(c, \ell, G) + \beta_d \cdot EV$\;
        \If{$val > V^{best}$}{
          $V^{best} \leftarrow val$; \quad $c^* \leftarrow c$; \quad $\ell^* \leftarrow \ell$\;
        }
      }
    }
    $V^n(a_j, z_s) \leftarrow V^{best}$\;
  }
  \If{$\|V^n - V^{n-1}\|_\infty < 10^{-6}$}{
    \textbf{break}\;
  }
}
Interpolate $(c^*, \ell^*)$ at each agent's actual $(a_{it}, z_{it})$\;
\end{algorithm}

\begin{remark}[Homogeneous Preference Speedup]
When all agents share identical utility parameters (auto-detected at runtime), the VFI in \Cref{alg:vfi} is solved once, producing policy grids $c^*(a, z)$ and $\ell^*(a, z)$. Each agent's decision is then obtained by interpolation at their state $(a_{it}, z_{it})$, reducing the computational cost from $O(N)$ VFI solves to one.
\end{remark}

\begin{remark}[NumPy Vectorization and Caching]
The VFI implementation uses NumPy vectorized operations throughout, replacing nested Python loops with array computations for approximately 100$\times$ speedup. Budget arrays, utility evaluations, and expected value interpolation are computed as batch array operations. Additionally, VFI results are cached with keys based on rounded market parameters (wage, interest rate, public goods, transfer, tax rate proxy, and utility weights). When prices are nearly identical across consecutive periods, the cache avoids redundant VFI solves entirely.
\end{remark}


\subsection{Endogenous Grid Method (EGM) Solver}\label{sec:egm}

As an alternative to VFI, the household problem can be solved via the Endogenous Grid Method (EGM) of \citet{carroll2006}. EGM avoids the inner maximization loop of VFI by inverting the Euler equation on an exogenous savings grid and recovering the endogenous current asset level. This yields substantial computational gains, particularly when combined with the Fella (2014) upper envelope for handling non-monotonicities near kinks in the budget constraint.

\subsubsection{Asset and Leisure Grids}

The EGM uses a 200-point exponentially-spaced savings grid on $[\underline{a},\, a_{max}]$ with $a_{max} = 1000$:
\begin{equation}\label{eq:egm_grid}
  a'_j = \underline{a} + (a_{max} - \underline{a}) \cdot \left(\frac{j}{N_a - 1}\right)^2, \qquad j = 0, 1, \ldots, N_a - 1, \quad N_a = 200.
\end{equation}
The quadratic spacing concentrates grid points near the borrowing constraint. A uniform leisure grid of $N_\ell = 51$ points is used for the intratemporal first-order condition.

\subsubsection{Euler Equation Inversion}

For each point on the exogenous savings grid $a'_j$ and each productivity state $z_s$, the EGM proceeds in five steps:

\begin{enumerate}[label=\textit{Step \arabic*.}]
  \item \textbf{Expected marginal utility.} Compute the right-hand side of the Euler equation:
  \begin{equation}\label{eq:egm_rhs}
    RHS(a'_j, z_s) = \beta_d (1+r) \sum_{s'} \Pi(s, s') \cdot u_c\!\left(c^{(n-1)}(a'_j, z_{s'}),\, \ell^{(n-1)}(a'_j, z_{s'}),\, G\right),
  \end{equation}
  where $u_c = \alpha_u \cdot c^{\alpha_u - 1} \cdot \ell^{\beta_u} \cdot G^{\gamma_u}$ is the marginal utility of consumption and superscript $(n-1)$ denotes the previous iteration's policy.

  \item \textbf{Euler inversion.} Invert the Euler equation to obtain consumption. Using the intratemporal FOC for leisure, $\ell^* = (\beta_u \cdot c) / (\alpha_u \cdot w \cdot z)$, substitute into the marginal utility condition and solve:
  \begin{equation}\label{eq:egm_inversion}
    c(a'_j, z_s) = \left(\frac{RHS}{\alpha_u \cdot \left(\frac{\beta_u}{\alpha_u \cdot w \cdot z_s}\right)^{\beta_u} \cdot G^{\gamma_u}}\right)^{1/(\alpha_u + \beta_u - 1)}.
  \end{equation}
  The intratemporal FOC then gives the associated leisure.

  \item \textbf{Endogenous grid.} Back out the current asset level from the budget constraint:
  \begin{equation}\label{eq:egm_endo}
    a(a'_j, z_s) = \frac{a'_j + c + T(wz_s(1-\ell)) - Tr - wz_s(1-\ell)}{1 + r}.
  \end{equation}
  This produces the \emph{endogenous} grid: the set of current asset levels that are consistent with saving $a'_j$ optimally.

  \item \textbf{Upper envelope.} The endogenous grid may be non-monotone near kinks in the budget constraint. Following \citet{fella2014}, we apply the upper envelope: sort the endogenous grid, remove dominated segments (keeping the higher-consumption branch at each crossing), and interpolate the cleaned policy back to the exogenous grid via monotone piecewise linear interpolation.

  \item \textbf{Constrained region.} For asset levels below the lowest point on the endogenous grid, the borrowing constraint binds: $a' = \underline{a}$. The agent consumes all available resources above $\underline{a}$, and leisure is recomputed from the intratemporal FOC at the constrained consumption level.
\end{enumerate}

\textbf{Convergence.} The EGM iterates until $\|c^{(n)} - c^{(n-1)}\|_\infty < 10^{-8}$, or a maximum of 500 iterations. An Euler equation residual check verifies accuracy: $\max |1 - c_{Euler}/c_{policy}| < 10^{-6}$ over unconstrained grid points.

\begin{remark}[Value Function Recovery]
After the EGM converges on the policy functions $c^*(a, z)$ and $\ell^*(a, z)$, the associated value function $V(a, z)$ is recovered by iterating the Bellman operator $V = u(c^*, \ell^*, G) + \lambda \cdot v(\calC; \theta) + \beta_d \, \E[V(a', z')]$ using the fixed policy until convergence. Here $\lambda \cdot v(\calC; \theta)$ is the political flow bonus described in \Cref{sec:political_utility}. The recovered value function is used by the entrepreneurial solver for occupational choice and by the political utility module for proposal evaluation.
\end{remark}


\subsection{Analytical Market Clearing}\label{sec:market_clearing}

Equilibrium prices $(w_t, r_t)$ are computed analytically from the Cobb-Douglas first-order conditions, yielding zero clearing error by construction.

Aggregate factor supplies are computed as:
\begin{equation}\label{eq:factor_supplies}
  L^s = \sum_{i: o_{it} = W} z_{it} \cdot \hat{h}_{it}, \qquad K^s = \max\!\left\{\sum_{i=1}^{N} a_{it},\; \sum_{f \in \calF_t} K_f\right\},
\end{equation}
where $\hat{h}_{it} = (1 - \ell_{it})$ if the household has a current labor supply decision, and $\hat{h}_{it} = 0.5$ as a fallback for stale labor supply values (i.e., before the first VFI solve). A floor guard of $10^{-8}$ is applied to both factor supplies to prevent division by zero.

With Cobb-Douglas production $Y = A_t (K^s)^\alpha (L^s)^{1-\alpha}$, the competitive equilibrium prices are:
\begin{equation}\label{eq:analytical_prices}
  w = (1-\alpha) \frac{Y}{L^s}, \qquad r = \alpha \frac{Y}{K^s} - \delta,
\end{equation}
subject to the floor guards $w \geq 10^{-6}$ and $r \geq -0.99$.

\begin{remark}[Zero Clearing Error]
Because prices are derived directly from the aggregate production function and factor supplies, this analytical approach produces identically zero market clearing error ($\epsilon = 0$) each period. This eliminates the convergence uncertainty of iterative methods.
\end{remark}

\begin{remark}[Historical Note]
An earlier version of the model used a tatonnement (price adjustment) algorithm that iteratively updated prices to equate firm-level factor demands with aggregate supply. The analytical approach dominates: it is faster, simpler, and exact. Tatonnement remains available as a conceptual alternative for models with non-Cobb-Douglas production where closed-form solutions do not exist.
\end{remark}


\subsection{Walrasian Market Clearing}\label{sec:walrasian}

When heterogeneous firms are present, the model supports Walrasian equilibrium via bisection on excess labor demand. Unlike the analytical approach, which uses a representative firm, Walrasian clearing respects firm-level heterogeneity in TFP and capital.

\subsubsection{Firm Factor Demands}

Each firm $f$ with TFP $A_f$ and capital $K_f$ has optimal labor demand from the Cobb-Douglas FOC:
\begin{equation}\label{eq:firm_labor_demand}
  L_f^*(w) = \left(\frac{(1-\alpha) \cdot A_f \cdot K_f^{\alpha}}{w}\right)^{1/\alpha}.
\end{equation}
Aggregate labor demand $L^d(w) = \sum_f L_f^*(w)$ is strictly decreasing in $w$, guaranteeing a unique equilibrium wage.

\subsubsection{Bisection Algorithm}

The equilibrium wage $w^*$ satisfies the excess labor demand condition:
\begin{equation}\label{eq:excess_labor}
  ELD(w) = L^d(w) - L^s = 0,
\end{equation}
where $L^s = \sum_{i: o_{it} = W} z_{it} (1 - \ell_{it})$ is aggregate effective labor supply from workers (excluding entrepreneurs, whose labor is embedded in firm production). The bisection algorithm:
\begin{enumerate}
  \item Bracket: find $w_{lo}$ with $ELD > 0$ and $w_{hi}$ with $ELD < 0$ (doubling $w_{hi}$ up to 50 times if needed).
  \item Bisect: iterate $w_{mid} = (w_{lo} + w_{hi})/2$, updating the bracket until $|ELD(w_{mid})| < 10^{-8}$ or 100 iterations.
\end{enumerate}

Given $w^*$, aggregate output is $Y = \sum_f A_f K_f^{\alpha} (L_f^*)^{1-\alpha}$, and the interest rate follows from the aggregate marginal product of capital: $r = \alpha Y / K_{total} - \delta$, where $K_{total} = \sum_f K_f$.

\begin{remark}[Fallback]
When no firms are present (e.g., in early periods before entrepreneurial entry), Walrasian clearing falls back to the representative-firm analytical formula of \Cref{sec:market_clearing}.
\end{remark}


\subsection{Entrepreneurial Decision Solver}\label{sec:entre_solver}

The entrepreneurial solver implements the utility-based occupational choice described in \Cref{sec:entry_exit}. The algorithm proceeds in four stages:
\begin{enumerate}
  \item \textbf{Stationary distribution.} At initialization, compute the stationary distribution $\pi^*$ of the ability Markov chain via power iteration (1000 iterations, tolerance $10^{-12}$). The expected stationary ability is $\bar{e} = \sum_j \pi^*_j \cdot e_j$.
  \item \textbf{Firm value.} For a given capital $K$, evaluate $V(e, K)$ using the closed-form formula \eqref{eq:firm_value}: current-period profit at actual ability plus the discounted perpetuity of stationary expected profit.
  \item \textbf{Optimal capital.} Find $K^*$ via grid search over 20 equally-spaced points in $[\underline{K}, a_{it} - F]$, maximizing $V(e_i, K)$.
  \item \textbf{Utility comparison.} Compute worker lifetime value $V^W$ \eqref{eq:worker_value} and entrepreneur lifetime value $V^E$ \eqref{eq:entre_value}, including the entry cost penalty $\Phi(F)$ for new entrants. The agent enters entrepreneurship if $V^E > V^W$ and $a_{it} \geq \underline{K} + F$; an existing entrepreneur exits if $V^W > V^E$.
\end{enumerate}

The solver also receives the current per-capita public goods level $G$ as input, since $G$ enters the utility function used for the occupational comparison.


\subsection{Bellman-Based Occupational Choice}\label{sec:bellman_occ}

When the Bellman occupational choice mode is active, the perpetuity-based firm value is replaced by a proper Bellman equation. The firm value $V^F(e, K)$ for all ability grid points is computed by solving:
\begin{equation}\label{eq:firm_bellman}
  V^F = (I - \beta_f \Pi_e)^{-1} \cdot \boldsymbol{\pi}(K),
\end{equation}
where $I$ is the $n_e \times n_e$ identity matrix, $\Pi_e$ is the ability transition matrix, $\beta_f$ is the firm discount factor, and $\boldsymbol{\pi}(K) = (\pi(e_1, K), \ldots, \pi(e_{n_e}, K))^\top$ is the vector of single-period profits at each ability grid point. The matrix $(I - \beta_f \Pi_e)^{-1}$ is pre-computed once at initialization and cached, making repeated firm value queries a single matrix-vector multiply. A condition number check ($< 10^{12}$) guards against near-singularity, falling back to the perpetuity formula if violated.

\subsubsection{Worker Value from VFI}

In Bellman mode, the worker value $V^W(a_i, z_i)$ is read directly from the VFI- or EGM-solved value function via linear interpolation on the asset grid at the household's actual state $(a_{it}, z_{it})$, rather than using the perpetuity approximation. This ensures exact consistency between the household optimization and the occupational choice comparison.

\subsubsection{Entrepreneur Leisure from FOC}

In Bellman mode, the entrepreneur's leisure is derived from the intratemporal FOC rather than hardcoded:
\begin{equation}\label{eq:entre_foc_leisure}
  \ell^*_E = \frac{\beta_u}{\alpha_u + \beta_u},
\end{equation}
which is the Cobb-Douglas time allocation that maximizes utility conditional on the profit stream, independent of prices. The continuation value uses $\max(V^W, V^F_{cont})$ to capture the option value of switching occupations each period.

\subsubsection{Golden-Section Search for Capital}

Optimal capital is found via golden-section search over $[K_{min}, a_{it} - F]$, evaluating the Bellman firm value at each candidate. The algorithm guarantees 50+ function evaluations for robust convergence, replacing the 20-point grid search used in legacy mode.


\subsection{Period Lifecycle}\label{sec:lifecycle}

Each period $t$ executes nine steps in a fixed order, guaranteeing determinism.

\begin{algorithm}[t]
\caption{Period Lifecycle (9-Step Algorithm)}\label{alg:lifecycle}
\KwIn{Previous period state $(t-1)$, configuration}
\KwOut{Updated period state $t$}
\textbf{Step 1: Draw shocks.} Transition productivity and ability via Markov chains; draw aggregate TFP from AR(1); optionally draw preference shocks.\;
\textbf{Step 2: Clear markets.} Find equilibrium $(w_t, r_t)$ via analytical Cobb-Douglas FOCs (\Cref{sec:market_clearing}).\;
\textbf{Step 3: Collect decisions.} Solve household Bellman via VFI (\Cref{alg:vfi}); solve utility-based entrepreneurial entry/exit (\Cref{sec:entre_solver}), receiving per-capita public goods $G$ for occupational comparison.\;
\textbf{Step 4: Validate constraints.} Project decisions onto feasible set: clamp $\ell \in [0,1]$, $c \leq budget - \underline{a}$, $a' \geq \underline{a}$.\;
\textbf{Step 5: Execute production.} Compute firm outputs via \eqref{eq:firm_production}; distribute income via \eqref{eq:firm_profit}; apply R\&D shocks via \eqref{eq:rd_prob}; create/liquidate firms.\;
\textbf{Step 6: Enforce constitution.} Apply tax rules, compute public goods, distribute transfers, apply regulations.\;
\textbf{Step 7: Process governance.} Collect proposals, validate, vote, amend constitution (at proposal-interval periods). In benchmark mode, uses rule-based governance with dissatisfaction-driven proposals and value-weighted voting (\Cref{sec:benchmark}).\;
\textbf{Step 8: Update states.} Apply DSGE-HA budget constraint: $a' = (1+r)a + y + \pi - c - T + Tr$; compute realized utility $u(c, \ell, G)$.\;
\textbf{Step 9: Observe.} Compute Gini, Pareto efficiency, aggregates (Y, C, I), wealth quantiles, welfare, firm stats; append to history.\;
\end{algorithm}

\begin{remark}[State Immutability]
Steps 1--9 operate on deep copies of the previous period's state. Agents never mutate global state directly; they return decisions that are validated and applied atomically in Step~8. This design prevents race conditions and ensures determinism (PROP-001).
\end{remark}


% =============================================================================
\section{Distribution Tracking via the Kolmogorov Forward Equation}\label{sec:distribution}
% =============================================================================

When the distribution mode is set to \texttt{kfe}, the model tracks the cross-sectional joint distribution of households over assets and productivity states on a discrete grid, rather than simulating individual agents. This approach, based on the Kolmogorov Forward Equation (KFE), is standard in the heterogeneous-agent literature for computing stationary distributions and aggregates.

\subsection{Distribution Object}

The distribution is a probability mass array $\mu(a_j, z_s)$ of dimension $N_a \times N_z$, where $\mu(a_j, z_s)$ represents the fraction of agents with assets $a_j$ and productivity $z_s$. The array satisfies $\sum_{j,s} \mu(a_j, z_s) = 1$ at all times (mass conservation).

\subsection{Young (2010) Lottery Allocation}\label{sec:young_lottery}

The KFE is advanced one period using the lottery allocation method of \citet{young2010}. For each grid cell $(a_j, z_s)$ with mass $\mu(a_j, z_s)$:

\begin{enumerate}
  \item \textbf{Policy lookup.} The savings policy gives $a' = g(a_j, z_s)$, the next-period asset level.
  \item \textbf{Grid bracket.} Find the bracket $a_{lo} \leq a' \leq a_{hi}$ on the asset grid.
  \item \textbf{Lottery weight.} Compute the interpolation weight for the upper bracket:
  \begin{equation}\label{eq:lottery_weight}
    \omega = \frac{a' - a_{lo}}{a_{hi} - a_{lo}}.
  \end{equation}
  \item \textbf{Mass allocation.} For each target productivity state $z_{s'}$, allocate mass using both the lottery weight and the Markov transition probability:
  \begin{align}
    \mu'(a_{hi}, z_{s'}) &\mathrel{+}= \omega \cdot \Pi(z_s, z_{s'}) \cdot \mu(a_j, z_s), \label{eq:kfe_hi} \\
    \mu'(a_{lo}, z_{s'}) &\mathrel{+}= (1 - \omega) \cdot \Pi(z_s, z_{s'}) \cdot \mu(a_j, z_s). \label{eq:kfe_lo}
  \end{align}
\end{enumerate}

The inner loops are vectorized using \texttt{np.add.at} for scatter-add operations, and a mass conservation check with renormalization is applied after each step (tolerance $10^{-12}$).

\subsection{Stationary Distribution}

The stationary (ergodic) distribution $\mu^*$ is computed by initializing with a uniform distribution and iterating the KFE forward step until convergence:
\begin{equation}\label{eq:kfe_convergence}
  \|\mu^{(n+1)} - \mu^{(n)}\|_1 < \epsilon_{KFE}, \qquad \epsilon_{KFE} = 10^{-10},
\end{equation}
with a maximum of 10{,}000 iterations. The stationary distribution is used for computing long-run aggregates and calibration targets.

\subsection{Aggregate Statistics from the Distribution}

Given the distribution $\mu$, aggregate quantities are computed as weighted sums:
\begin{equation}\label{eq:kfe_aggregate}
  X = \sum_{j,s} \mu(a_j, z_s) \cdot f(a_j, z_s),
\end{equation}
where $f$ is any policy function (consumption, labor supply, savings). The Gini coefficient is computed from the marginal wealth distribution $\mu_a(a_j) = \sum_s \mu(a_j, z_s)$ using the trapezoidal Lorenz curve formula:
\begin{equation}\label{eq:kfe_gini}
  G_{KFE} = 1 - 2 \int_0^1 L(p)\, dp,
\end{equation}
where $L(p)$ is the Lorenz curve constructed from cumulative wealth shares. Additional distributional statistics---percentiles, top wealth shares, and mean wealth---are computed directly from $\mu_a$.

\subsection{Agent Sampling}

For compatibility with agent-level operations (e.g., governance and LLM interaction), individual agents can be sampled from the distribution via inverse CDF sampling on the flattened joint distribution, using the simulation's seeded RNG for determinism.


% =============================================================================
\section{Two-Asset Household Structure}\label{sec:two_asset}
% =============================================================================

When two-asset mode is active, the single-asset framework is extended to a two-asset model following \citet{kaplanmollviolante2018}. Households hold two types of assets: liquid government bonds $b$ and illiquid physical capital $k$.

\subsection{Portfolio Structure}

The household state is augmented from $(a, z, e, o)$ to $(b, k, z, e, o)$, where $b \in [b_{min}, \infty)$ is liquid bond holdings (which may be negative for unsecured credit when $b_{min} < 0$) and $k \in [0, \infty)$ is illiquid capital. Total wealth is $a = b + k$.

Each period, the household chooses a deposit $d = k' - k$ into the illiquid asset, subject to a convex adjustment cost:
\begin{equation}\label{eq:adj_cost}
  \chi(d, k) = \chi_0 |d| + \chi_1 \frac{d^2}{\max(k, 1)},
\end{equation}
where $\chi_0 \geq 0$ is the linear component and $\chi_1 \geq 0$ is the quadratic component. The denominator $\max(k, 1)$ prevents division by zero for agents with no illiquid holdings.

\subsection{Budget Constraint}

The two-asset budget constraint for a worker is:
\begin{equation}\label{eq:two_asset_budget}
  b' = (1 + r^b) b + wz(1-\ell) - c - T(y) + Tr - d - \chi(d, k),
\end{equation}
where $r^b$ is the bond rate from the Fisher equation and $d + \chi(d, k)$ is the total cost of adjusting illiquid holdings.

\subsection{``Wealthy Hand-to-Mouth'' Agents}

A key feature of the two-asset model is the emergence of ``wealthy hand-to-mouth'' agents: households with substantial illiquid wealth (e.g., housing or business capital) but little liquid wealth. These agents have high marginal propensities to consume (MPCs) out of transitory income despite being wealthy in total, because accessing their illiquid assets is costly. The distribution of agents across the $(b, k)$ space---and in particular the mass of agents near $b \approx b_{min}$ with $k > 0$---is a key determinant of the economy's response to fiscal and monetary policy shocks.

\subsection{Nested EGM}

In two-asset mode, the EGM solver operates on a three-dimensional state space $(b, k, z)$. The solution uses a nested approach: for each illiquid capital level on a 100-point grid, the one-dimensional EGM for the liquid asset is solved conditional on the deposit choice. A 30-point deposit grid is used for the outer portfolio decision. The overall structure follows the nested EGM approach for two-asset models.


% =============================================================================
\section{Government Sector and Bond Market}\label{sec:government}
% =============================================================================

The government sector issues bonds to finance spending and transfers in excess of tax revenue. A fiscal rule ensures long-run debt sustainability.

\subsection{Government Budget Constraint}

The government's flow budget constraint is:
\begin{equation}\label{eq:govt_budget}
  B_{t+1} = (1 + r^b_t) B_t + G_t + Tr_t - T_t,
\end{equation}
where $B_t$ is outstanding government debt (bonds), $r^b_t$ is the bond interest rate, $G_t$ is public goods spending, $Tr_t$ is total transfers, and $T_t$ is total tax revenue. When tax revenue exceeds spending plus debt service, debt decreases; otherwise it increases.

\subsection{Bond Market Clearing}\label{sec:bond_clearing}

The bond rate $r^b_t$ is determined by clearing the bond market. Aggregate household bond demand is modeled as a logistic function of the spread between the bond rate and the capital rate:
\begin{equation}\label{eq:bond_demand}
  B^d(r^b) = \sum_{i=1}^{N} a_{it} \cdot \sigma(s \cdot (r^b - r)),
\end{equation}
where $\sigma(\cdot)$ is the logistic function, $s = 5$ is the sensitivity parameter, and $r$ is the capital market interest rate (the opportunity cost of holding bonds).

The equilibrium bond rate $r^{b*}$ satisfies $B^d(r^{b*}) = B_t$ and is found via bisection on $r^b \in [-0.5, 1.0]$ with tolerance $10^{-8}$ and a maximum of 100 iterations.

\subsection{Fiscal Rule}\label{sec:fiscal_rule}

To prevent explosive debt paths, a fiscal rule automatically adjusts tax rates when the debt-to-GDP ratio exceeds a threshold:
\begin{equation}\label{eq:fiscal_rule}
  \Delta \tau_t =
  \begin{cases}
    \tau_{adj} & \text{if } B_t / Y_t > \bar{B}, \\
    0 & \text{otherwise},
  \end{cases}
\end{equation}
where $\bar{B}$ is the debt-to-GDP threshold (default 1.5) and $\tau_{adj}$ is the tax rate increment (default 0.01). A safety cap flags critical warnings when $B_t / Y_t > 5$, and governance changes are frozen to prevent further fiscal deterioration.


% =============================================================================
\section{Nominal Rigidities}\label{sec:nominal}
% =============================================================================

When nominal rigidities are enabled, the model incorporates sticky prices, a monetary policy rule, and the Fisher equation linking nominal and real rates.

\subsection{Rotemberg Sticky Prices}\label{sec:rotemberg}

Price adjustment follows \citet{rotemberg1982}, with firms facing a quadratic cost of changing prices. In the symmetric equilibrium, the New Keynesian Phillips Curve (NKPC) is:
\begin{equation}\label{eq:nkpc}
  \phi_p \pi_t (\pi_t - \bar{\pi}) = (1 - \varepsilon) + \varepsilon \cdot mc_t + \beta_d \phi_p \E_t\!\left[\pi_{t+1}(\pi_{t+1} - \bar{\pi}) \frac{Y_{t+1}}{Y_t}\right],
\end{equation}
where $\phi_p$ is the Rotemberg price adjustment cost, $\pi_t$ is inflation, $\bar{\pi}$ is the inflation target, $\varepsilon$ is the elasticity of substitution between varieties, and $mc_t = w_t / A_t$ is the real marginal cost.

This is a quadratic equation in $\pi_t$. Given expectations $\E_t[\pi_{t+1}]$ and $\E_t[Y_{t+1}]$, the NKPC is solved for the inflation rate $\pi_t$ that is closest to the target $\bar{\pi}$ (the stable root).

\subsection{Taylor Rule}\label{sec:taylor}

Monetary policy follows a Taylor rule \citep{taylor1993}:
\begin{equation}\label{eq:taylor}
  i_t = r^{nat} + \phi_\pi (\pi_t - \bar{\pi}) + \phi_y \cdot gap_t,
\end{equation}
where $i_t$ is the nominal interest rate, $r^{nat}$ is the natural (real) interest rate (proxied by the capital market rate), $\phi_\pi > 1$ is the inflation coefficient satisfying the Taylor principle, $\phi_y \geq 0$ is the output gap coefficient, and $gap_t = (Y_t - \bar{Y})/\bar{Y}$ is the output gap relative to steady-state output $\bar{Y}$.

\subsection{Fisher Equation}

The real bond rate is determined by the Fisher equation:
\begin{equation}\label{eq:fisher}
  r^b_t = \frac{1 + i_t}{1 + \E_t[\pi_{t+1}]} - 1,
\end{equation}
linking the nominal rate from the Taylor rule and inflation expectations to the real return on government bonds that enters household budget constraints.

\subsection{Nominal State and Convergence}

The nominal block tracks a \texttt{NominalState} object each period consisting of the price level $P_t$, inflation $\pi_t$, nominal rate $i_t$, real bond rate $r^b_t$, expected inflation $\E_t[\pi_{t+1}]$, and marginal cost $mc_t$.

The three equations---NKPC, Taylor rule, and Fisher equation---form a simultaneous system. It is solved via an outer iteration:
\begin{enumerate}
  \item Initialize expected inflation using adaptive expectations: $\E[\pi_{t+1}] = 0.9 \cdot \pi_{t-1}$.
  \item Compute marginal cost $mc_t = w_t / A_t$.
  \item Solve Taylor rule $\to$ $i_t$, Fisher equation $\to$ $r^b_t$, NKPC $\to$ $\pi_t$.
  \item Update $\E[\pi_{t+1}] = 0.9 \cdot \pi_t$ and iterate until $|\pi_t^{(k)} - \pi_t^{(k-1)}| < 10^{-6}$ or 50 iterations.
\end{enumerate}

The price level evolves as $P_t = P_{t-1}(1 + \pi_t)$.


% =============================================================================
\section{Microfounded Political Utility}\label{sec:political_utility}
% =============================================================================

A key contribution is the microfoundation of political preferences within the household's Bellman equation. Rather than treating political decisions as ad hoc or purely ideological, we embed a political utility function in the agent's value function, yielding structurally consistent political behavior derived from economic self-interest and ideology.

\subsection{Political Utility Function}

Each agent $i$ has a political utility function:
\begin{equation}\label{eq:political_utility}
  v(\calC; \theta_i) = \theta_i^{eq} \cdot f^{eq}(\calC) + \theta_i^{lib} \cdot f^{lib}(\calC),
\end{equation}
where $\theta_i^{eq}$ and $\theta_i^{lib}$ are the agent's value vector weights (from the ideological position described in \Cref{sec:preferences}), and:
\begin{align}
  f^{eq}(\calC) &= -\text{Gini}(\calC), \label{eq:f_eq} \\
  f^{lib}(\calC) &= -\tau(\calC), \label{eq:f_lib}
\end{align}
where $\text{Gini}(\calC)$ is the Gini coefficient under the current constitution (or a structural proxy when simulation-level Gini is unavailable) and $\tau(\calC)$ is the effective tax rate. Equality-oriented agents prefer lower Gini; liberty-oriented agents prefer lower taxes.

\subsection{Bellman Integration}

The political utility enters the household's Bellman equation as a constant flow bonus:
\begin{equation}\label{eq:bellman_political}
  V(a, z) = \max_{c, \ell} \left\{u(c, \ell, G) + \lambda \cdot v(\calC; \theta) + \beta_d \, \E[V(a', z')]\right\},
\end{equation}
where $\lambda \geq 0$ is the weight on political utility (default 0.05). Since the constitution $\calC$ is fixed between governance periods, $v(\calC; \theta)$ is a constant that does not affect the consumption-savings-labor first-order conditions but does affect the \emph{level} of the value function. This level matters for:
\begin{itemize}
  \item Occupational choice (comparing $V^W$ vs.\ $V^E$ under different constitutions);
  \item Proposal evaluation (comparing $V$ under current vs.\ proposed rules);
  \item Voting decisions (support proposals that increase $V$).
\end{itemize}

\subsection{Proposal Evaluation via Value Function Comparison}

When evaluating a proposed constitutional change $\calC' \neq \calC$, the welfare gain for agent $i$ is:
\begin{equation}\label{eq:proposal_eval}
  \Delta V_i = \underbrace{\frac{\lambda}{1 - \beta_d}\left[v(\calC'; \theta_i) - v(\calC; \theta_i)\right]}_{\text{political component}} + \underbrace{\left(-\frac{\Delta\tau \cdot y_i}{1 - \beta_d}\right)}_{\text{economic component}},
\end{equation}
where the political component is the capitalized (perpetuity) value of the political utility change, and the economic component approximates the disposable income effect of the tax rate change $\Delta\tau = \tau(\calC') - \tau(\calC)$.

Agent $i$ supports the proposal if $\Delta V_i > 0$. When \texttt{pure\_bellman\_politics} is active, this comparison replaces the LLM for all political decisions, yielding fully microfounded governance dynamics.

\subsection{LLM Integration with Bellman Context}

When the LLM is used for political decisions, the Bellman-derived preference is provided as structured context: the recommendation (support/oppose), the magnitude of $\Delta V$, and the decomposition into political and economic components. The LLM may override the Bellman recommendation for strategic or coalition reasons; such deviations are logged for analysis.


% =============================================================================
\section{LLM-Augmented Political Decisions}\label{sec:llm}
% =============================================================================

\subsection{Motivation}\label{sec:llm_motivation}

The model distinguishes between two types of decisions. \emph{Economic decisions}---consumption, savings, labor supply---have well-defined objective functions (the Bellman equation) and can be solved numerically to arbitrary precision. \emph{Political decisions}---what tax rate to propose, whether to add a means-tested transfer, how to vote on a minimum wage---are fundamentally open-ended. The space of possible proposals is large, the evaluation criteria are multi-dimensional (personal utility, ideological values, societal welfare), and the strategic interactions are complex.

This asymmetry motivates our design: analytical optimization for economics, LLM-augmented bounded rationality for politics. The LLM acts not as an optimizer but as a model of deliberative decision-making under genuine uncertainty about institutional design.


\subsection{Architecture}\label{sec:llm_arch}

The LLM decision engine operates through a five-stage pipeline:

\begin{enumerate}
  \item \textbf{Context construction.} For each agent, build a structured JSON context containing: (a) personal state (wealth, productivity, role, utility parameters, value vector); (b) economic conditions (wage, interest rate, aggregate output); (c) current constitution (all rules and their parameters); (d) recent history (Gini trends, welfare changes, proposal outcomes); (e) a rule-type guide with examples and parameter descriptions.

  \item \textbf{Cache lookup.} Compute a hash of each agent's context. If an identical context was seen earlier in the simulation, reuse the cached response. This avoids redundant API calls when multiple agents face similar conditions.

  \item \textbf{Batched API calls.} Uncached contexts are batched (default batch size: 10) and sent to the LLM provider with structured output schemas. The LLM returns JSON conforming to the \texttt{PoliticalDecision} schema, which includes an optional proposal and a vote map.

  \item \textbf{Schema validation.} Each LLM response is validated against a Pydantic model. Responses that fail validation trigger the fallback path.

  \item \textbf{Fallback.} On API failure, timeout, or validation error, the agent defaults to a null political decision (no proposal, no votes), ensuring the simulation always progresses.
\end{enumerate}


\subsection{Information Structure}\label{sec:llm_info}

The political context provided to the LLM is richer than the economic context, reflecting the complexity of institutional design:

\begin{itemize}
  \item \textbf{Societal conditions:} Gini coefficient, mean and median wealth, wealth quantiles (p10, p25, p50, p75, p90), unemployment rate, number of active firms, aggregate output, social welfare.

  \item \textbf{Recent history with trends:} The last 3--5 observation entries, each annotated with trend descriptors (``Inequality decreased,'' ``Welfare improved,'' etc.) computed by comparing consecutive observations.

  \item \textbf{Recent proposal outcomes:} The last 20 vote outcomes, including whether proposals were passed, voted down, or rejected at validation, with vote counts. This enables the LLM to learn from prior political dynamics and avoid repeating failed approaches.

  \item \textbf{Rule-type guide:} A structured dictionary describing each of the eight rule types, their valid parameters, and example proposals. This serves as an instruction manual for the LLM's institutional design.
\end{itemize}


\subsection{Determinism and Reproducibility}\label{sec:llm_determinism}

Several design choices ensure reproducibility:
\begin{itemize}
  \item LLM temperature is set to $0$, producing deterministic outputs for identical inputs.
  \item All randomness in agent selection, proposal ordering, and tie-breaking flows through a single seeded RNG.
  \item The cache maps contexts to responses deterministically.
\end{itemize}


\subsection{Benchmark Mode}\label{sec:benchmark}

The model supports a \emph{benchmark mode} that disables LLM calls entirely. In benchmark mode:
\begin{itemize}
  \item Economic decisions are computed by the VFI solver.
  \item Entrepreneurial decisions are computed by the utility-based occupational choice solver.
  \item Governance uses rule-based decision-making via deterministic citizen logic.
\end{itemize}

Unlike earlier versions where the constitution remained static, benchmark mode now features a \emph{rule-based governance} system. Each agent's political behavior is determined algorithmically:

\begin{itemize}
  \item \textbf{Proposal gating.} Each agent skips proposal generation with probability 0.85, so that approximately 15\% of agents propose in each governance period.
  \item \textbf{Dissatisfaction-based proposals.} Proposing agents evaluate their dissatisfaction with current constitutional rules based on their value vector $(v^{eq}, v^{lib})$. For instance, an agent with high $v^{eq}$ is dissatisfied if taxes are low relative to their preferred rate $v^{eq} \times 0.5$. The agent proposes a modification to the rule with highest dissatisfaction, adjusting the parameter 30\% toward their preferred value.
  \item \textbf{Value-weighted voting.} Each agent scores proposals using a weighted combination:
  \begin{equation}\label{eq:benchmark_voting}
    \text{Score} = 0.6 \times \text{value\_alignment} + 0.4 \times \text{self\_interest},
  \end{equation}
  where value alignment measures how the proposal moves the rule toward the agent's ideological preference, and self-interest reflects the material impact (e.g., wealthier agents oppose higher taxes). The agent votes \emph{for} if $\text{Score} > 0.5$.
\end{itemize}

This enables controlled experiments that compare LLM-augmented governance against a well-defined rule-based baseline, rather than a static constitution.


% =============================================================================
\section{Calibration}\label{sec:calibration}
% =============================================================================

\Cref{tab:calibration} reports the baseline calibration. Parameter values are drawn from standard sources in the quantitative macro literature.

\begin{table}[tp]
\centering
\caption{Baseline Calibration}\label{tab:calibration}
\small
\begin{tabular}{@{}llcl@{}}
\toprule
\textbf{Parameter} & \textbf{Symbol} & \textbf{Value} & \textbf{Description} \\
\midrule
\multicolumn{4}{l}{\textit{Core}} \\
Number of agents & $N$ & 50 & Population size \\
Time horizon & $T$ & 100 & Number of periods \\
RNG seed & --- & 42 & Reproducibility \\
\midrule
\multicolumn{4}{l}{\textit{Production}} \\
Capital share & $\alpha$ & 0.33 & Cobb-Douglas \\
Depreciation rate & $\delta$ & 0.10 & Annual \\
Aggregate TFP (initial) & $A_0$ & 1.0 & Normalized \\
\midrule
\multicolumn{4}{l}{\textit{Household preferences}} \\
Consumption weight & $\alpha_u$ & 0.40 & Cobb-Douglas utility \\
Leisure weight & $\beta_u$ & 0.35 & Cobb-Douglas utility \\
Public goods weight & $\gamma_u$ & 0.25 & Cobb-Douglas utility \\
Discount factor & $\beta_d$ & 0.95 & Intertemporal \\
Borrowing constraint & $\underline{a}$ & 0.0 & Natural borrowing limit \\
Initial wealth mean & $\bar{a}_0$ & 100.0 & Log-normal approx \\
Initial wealth std & $\sigma_{a_0}$ & 30.0 & --- \\
\midrule
\multicolumn{4}{l}{\textit{Idiosyncratic shocks}} \\
Productivity persistence & $\rho_z$ & 0.90 & AR(1) \\
Productivity volatility & $\sigma_z$ & 0.20 & Innovation s.d. \\
Grid points & $N_z$ & 7 & Rouwenhorst \\
\midrule
\multicolumn{4}{l}{\textit{Entrepreneurial}} \\
Ability persistence & $\rho_e$ & 0.85 & AR(1) \\
Ability volatility & $\sigma_e$ & 0.30 & Innovation s.d. \\
Grid points & $N_e$ & 7 & Rouwenhorst \\
Entry cost & $F$ & 5.0 & Sunk cost \\
Minimum capital & $\underline{K}$ & 10.0 & --- \\
Firm discount factor & $\beta_f$ & 0.95 & Perpetuity discounting \\
Entrepreneur leisure & $\ell_E$ & 0.20 & Reduced leisure (effort) \\
\midrule
\multicolumn{4}{l}{\textit{Aggregate shocks}} \\
TFP persistence & $\rho_A$ & 0.95 & AR(1) \\
TFP volatility & $\sigma_A$ & 0.01 & Innovation s.d. \\
\midrule
\multicolumn{4}{l}{\textit{R\&D}} \\
Base success probability & $p_0$ & 0.10 & --- \\
TFP improvement mean & $\mu_{rd}$ & 0.05 & --- \\
TFP improvement std & $\sigma_{rd}$ & 0.02 & --- \\
\bottomrule
\end{tabular}
\end{table}

\begin{table}[tp]
\centering
\caption{Baseline Calibration (continued)}\label{tab:calibration2}
\small
\begin{tabular}{@{}llcl@{}}
\toprule
\textbf{Parameter} & \textbf{Symbol} & \textbf{Value} & \textbf{Notes} \\
\midrule
\multicolumn{4}{l}{\textit{Market clearing}} \\
Method & --- & Analytical/Walrasian & Configurable \\
Walrasian bisection tol & --- & $10^{-8}$ & --- \\
Walrasian max iter & --- & 100 & --- \\
\midrule
\multicolumn{4}{l}{\textit{EGM solver}} \\
Asset grid points & $N_a$ & 200 & Exponential spacing \\
Leisure grid points & $N_\ell$ & 51 & Uniform \\
Convergence tolerance & $\epsilon_{EGM}$ & $10^{-6}$ & Sup-norm \\
\midrule
\multicolumn{4}{l}{\textit{VFI (fallback)}} \\
Asset grid points & $N_a$ & 100 & Exponential spacing \\
Leisure grid points & $N_\ell$ & 21 & Uniform \\
Max iterations & --- & 500 & --- \\
Convergence tolerance & $\epsilon_{VFI}$ & $10^{-6}$ & Sup-norm \\
\midrule
\multicolumn{4}{l}{\textit{Two-asset (optional)}} \\
Linear adj.\ cost & $\chi_0$ & 0.01 & $\geq 0$ \\
Quadratic adj.\ cost & $\chi_1$ & 0.005 & $\geq 0$ \\
Liquid borrowing limit & $b_{min}$ & 0.0 & --- \\
\midrule
\multicolumn{4}{l}{\textit{Government (optional)}} \\
Initial debt & $B_0$ & 0.0 & $\geq 0$ \\
Debt/GDP threshold & $\bar{B}$ & 1.5 & $> 0$ \\
Fiscal rule adj. & $\tau_{adj}$ & 0.01 & $(0, 1)$ \\
\midrule
\multicolumn{4}{l}{\textit{Nominal (optional)}} \\
Rotemberg cost & $\phi_p$ & 100 & $\geq 0$ \\
Taylor infl.\ coeff. & $\phi_\pi$ & 1.5 & $> 1$ \\
Taylor output coeff. & $\phi_y$ & 0.125 & $\geq 0$ \\
Inflation target & $\bar{\pi}$ & 0.02 & --- \\
Elasticity of sub. & $\varepsilon$ & 6.0 & $> 1$ \\
\midrule
\multicolumn{4}{l}{\textit{Political utility (optional)}} \\
Political weight & $\lambda$ & 0.05 & $[0, 1]$ \\
\midrule
\multicolumn{4}{l}{\textit{Governance}} \\
Default tax rate & $\tau_0$ & 0.10 & Initial flat tax (10\%) \\
Proposal interval & $K_{prop}$ & 5 & Periods \\
Observer interval & $K_{obs}$ & 5 & Periods \\
Proposal skip probability & --- & 0.85 & Benchmark mode gating \\
\midrule
\multicolumn{4}{l}{\textit{LLM}} \\
Temperature & --- & 0.0 & Deterministic \\
Batch size & --- & 10 & Agents per call \\
Cache enabled & --- & Yes & --- \\
\bottomrule
\end{tabular}
\end{table}

The capital share $\alpha = 0.33$ and depreciation rate $\delta = 0.10$ are standard in the macro literature. The productivity process parameters ($\rho_z = 0.9$, $\sigma_z = 0.2$) are in the range estimated by \citet{storesletten2004} for annual earnings. The Rouwenhorst discretization uses 7 grid points for both productivity and ability, which \citet{kopeckysuen2010} show provides adequate approximation for persistence above 0.8. The entrepreneurial ability process is somewhat more volatile ($\sigma_e = 0.3$) and less persistent ($\rho_e = 0.85$), reflecting the higher turnover in business outcomes relative to labor earnings.

The discount factor $\beta_d = 0.95$ and the Cobb-Douglas utility weights ($\alpha_u = 0.4$, $\beta_u = 0.35$, $\gamma_u = 0.25$) are within the range of estimated preferences in the literature on labor supply and public goods valuation. The EGM solver uses a 200-point exponential asset grid and 51-point leisure grid, providing substantially finer resolution than the VFI fallback.

The two-asset adjustment costs ($\chi_0 = 0.01$, $\chi_1 = 0.005$) follow \citet{kaplanmollviolante2018} in generating a realistic mass of ``wealthy hand-to-mouth'' agents. The Rotemberg price adjustment cost $\phi_p = 100$ and Taylor rule coefficients ($\phi_\pi = 1.5$, $\phi_y = 0.125$) are standard in the New Keynesian literature. The political utility weight $\lambda = 0.05$ is set conservatively so that political preferences matter for institutional choice but do not dominate economic decisions.


\subsection{Simulated Method of Moments (SMM) Calibration}\label{sec:smm}

For systematic parameter estimation, the model supports Simulated Method of Moments (SMM) calibration. The calibration procedure matches model-implied moments to empirical targets by minimizing a weighted sum of squared deviations.

\subsubsection{Calibration Targets}

The default calibration targets are:
\begin{itemize}
  \item Wealth Gini coefficient: 0.80 (consistent with U.S.\ SCF data).
  \item Entrepreneur share: 0.10 (fraction of population).
  \item Top 10\% wealth share: 0.70.
\end{itemize}
Additional targets for the two-asset extension include the median marginal propensity to consume (MPC) and the liquid-to-illiquid wealth ratio.

\subsubsection{SMM Objective}

The SMM objective function is:
\begin{equation}\label{eq:smm}
  Q(\theta) = \sum_{m \in \calM} \omega_m \left(\frac{\hat{m}(\theta) - m^{data}}{m^{data}}\right)^2,
\end{equation}
where $\theta$ is the parameter vector, $\hat{m}(\theta)$ is the model-implied moment, $m^{data}$ is the empirical target, and $\omega_m$ is the moment weight (default: equal weights). The relative deviation $(\hat{m} - m^{data})/m^{data}$ ensures scale-invariance across moments.

\subsubsection{Optimization}

The SMM objective is minimized via \texttt{scipy.optimize.minimize} using the Nelder-Mead (simplex) method, which is derivative-free and robust to the discontinuities that arise from the simulation. The calibrated parameter vector typically includes the discount factor $\beta_d$, productivity volatility $\sigma_z$, entry cost $F$, and borrowing limit $\underline{a}$. Parameters are projected onto feasible bounds at each evaluation (e.g., $\beta_d \in [0.80, 0.999]$, $\sigma_z \in [0.05, 1.0]$, $F \in [0.1, 50.0]$). Failed simulations are penalized with a large objective value ($10^{10}$) to steer the optimizer away from unstable regions.


% =============================================================================
\section{Discussion}\label{sec:discussion}
% =============================================================================

\subsection{Wealth Inequality Dynamics}

The model generates endogenous wealth inequality through three channels. First, idiosyncratic productivity shocks create earnings inequality among workers. Second, entrepreneurial returns amplify the right tail of the wealth distribution, as successful entrepreneurs earn profits in addition to labor income. Third, the constitutional feedback loop creates a political economy of inequality: if wealthier agents successfully block redistributive proposals (exploiting their economic position to influence governance outcomes), inequality may be self-reinforcing.

The Gini coefficient is computed at each observation interval as
\begin{equation}\label{eq:gini}
  G = \frac{2 \sum_{i=1}^{N} i \cdot a_{(i)} - (N+1) \sum_{i=1}^{N} a_{(i)}}{N \sum_{i=1}^{N} a_{(i)}},
\end{equation}
where $a_{(1)} \leq a_{(2)} \leq \cdots \leq a_{(N)}$ are ordered wealth levels.


\subsection{Occupational Choice Mechanisms}

The entry/exit dynamics are driven by the interaction of three factors:
\begin{enumerate}
  \item \textbf{Ability:} High-ability agents have higher expected firm values, making entry more attractive.
  \item \textbf{Wealth:} The minimum capital constraint $\underline{K}$ and entry cost $F$ create a wealth barrier to entrepreneurship. Agents must satisfy $a_{it} \geq \underline{K} + F$ to enter.
  \item \textbf{Prices:} Factor prices $(w, r)$ affect both the opportunity cost of entrepreneurship (the forgone return $a \cdot r$ on savings) and the profitability of firms (through labor costs and capital costs).
\end{enumerate}

The interaction of these factors generates rich dynamics. Rising wages reduce firm profitability, potentially triggering exit. A constitutional increase in tax rates reduces after-tax profits but may also reduce the opportunity cost of capital if it lowers interest rates through general equilibrium effects.


\subsection{Constitutional Evolution}

In LLM-augmented mode, the constitution evolves as agents observe economic outcomes and propose institutional changes. Several patterns emerge:
\begin{itemize}
  \item \textbf{Tax ratchet:} If inequality is high, equality-preferring agents may propose higher taxes. If the tax rate becomes too high, liberty-preferring agents propose reductions. This creates oscillatory dynamics around an interior equilibrium.
  \item \textbf{Transfer design:} Agents may initially propose equal-share transfers, then shift to means-tested transfers as they observe that flat transfers have limited equalizing effect.
  \item \textbf{Voting rule changes:} Constitutional amendments to the voting threshold itself create meta-governance dynamics---agents may try to entrench favorable rules by raising the threshold required to change them.
\end{itemize}


\subsection{LLM vs.\ Benchmark Comparison}

Comparing LLM-mode and benchmark-mode simulations isolates the effect of LLM-augmented deliberation versus rule-based governance. In benchmark mode, the constitution evolves through rule-based governance with dissatisfaction-driven proposals and value-weighted voting (\Cref{sec:benchmark}), starting from the default constitution (10\% flat tax, equal-share transfers, majority voting). In LLM mode, constitutional changes are driven by LLM-generated proposals informed by rich economic context, potentially producing more nuanced institutional designs. Both modes feature endogenous governance, but they differ in the sophistication of the political reasoning.

The welfare comparison is measured by cumulative discounted social welfare:
\begin{equation}\label{eq:cumulative_welfare}
  W_T = \sum_{t=0}^{T} \bar{\beta}^t \sum_{i=1}^{N} u(c_{it}, \ell_{it}, G_t),
\end{equation}
where $\bar{\beta}$ is the average discount factor across agents.


\subsection{Pareto Efficiency}

The Pareto efficiency score estimates the fraction of agent pairs for which no improving transfer exists. For each ordered pair $(i, j)$ with $a_i \geq a_j$, we test whether a small transfer $\Delta = 1.0$ from $i$ to $j$ would improve $j$'s utility without reducing $i$'s:
\begin{equation}\label{eq:pareto_test}
  u_j(c_j + \Delta) > u_j(c_j) \quad \text{and} \quad u_i(c_i - \Delta) \geq u_i(c_i).
\end{equation}
The Pareto score is $1$ minus the fraction of pairs where such an improvement exists.


% =============================================================================
\section{Conclusion}\label{sec:conclusion}
% =============================================================================

We have developed a Heterogeneous Agent New Keynesian (HANK) model in which governance institutions are fully endogenous. Agents with heterogeneous productivity, ability, and preferences simultaneously make economic decisions (consumption, savings, labor, entrepreneurship, portfolio allocation) and political decisions (proposals, votes on constitutional rules). The economic and political layers are tightly coupled: constitutional rules determine fiscal policy, which enters budget constraints and affects equilibrium prices, which in turn shape the political preferences that generate new proposals.

The computational framework combines the Endogenous Grid Method for household optimization with Walrasian market clearing, Kolmogorov Forward Equation distribution tracking, a government sector with bond market, Rotemberg nominal rigidities with a Taylor rule, microfounded political utility in the Bellman equation, and LLM-augmented decision-making for political deliberation. This design reflects a genuine asymmetry in the underlying problems: economic optimization has well-defined structure (the Bellman equation), while institutional design is genuinely open-ended.

Several extensions merit future investigation. First, the model could incorporate a richer set of agent interactions, including coalition formation and bargaining. Second, the LLM's ``preferences'' could be calibrated against survey data on political attitudes. Third, the constitutional framework could be extended to include judicial review and enforcement uncertainty. Fourth, the life-cycle/OLG structure could be added for richer demographic dynamics.


% =============================================================================
% Appendices
% =============================================================================
\appendix

\section{Rouwenhorst Method: Detailed Construction}\label{app:rouwenhorst}

We provide the full algorithm for the Rouwenhorst discretization of an AR(1) process $\log x_t = \rho \log x_{t-1} + \varepsilon_t$ with $\varepsilon_t \sim \mathcal{N}(0, \sigma^2)$.

\textbf{Input:} Persistence $\rho \in (0, 1)$, innovation volatility $\sigma > 0$, number of grid points $n \geq 2$.

\textbf{Output:} Grid values $\calS = \{s_1, \ldots, s_n\}$ in levels, transition matrix $\Pi \in \R^{n \times n}$.

\begin{enumerate}
  \item Unconditional standard deviation: $\sigma_x = \sigma / \sqrt{1 - \rho^2}$.
  \item Grid bounds in logs: $z_{max} = \sigma_x \sqrt{n-1}$.
  \item Log grid: $z_j = -z_{max} + (j-1) \cdot 2z_{max}/(n-1)$ for $j = 1, \ldots, n$.
  \item Transition parameter: $p = (1 + \rho)/2$.
  \item Base matrix ($n = 2$):
  \[
    \Pi_2 = \begin{pmatrix} p & 1-p \\ 1-p & p \end{pmatrix}.
  \]
  \item Recursive construction. For $m = 3, \ldots, n$, define the $(m \times m)$ matrix $\widetilde{\Pi}_m$ by embedding $\Pi_{m-1}$ in four quadrants (see \eqref{eq:rouwenhorst_recursion}), then normalize rows to sum to 1.
  \item Level grid: $s_j = \exp(z_j)$ for $j = 1, \ldots, n$.
\end{enumerate}

The stationary distribution $\pi^*$ satisfying $\pi^* = \pi^* \Pi$ is computed by initializing $\pi^{(0)} = (1/n, \ldots, 1/n)$ and iterating $\pi^{(k+1)} = \pi^{(k)} \Pi$ until convergence.


\section{VFI Convergence Properties}\label{app:vfi}

The VFI algorithm inherits convergence guarantees from the contraction mapping theorem. Define the Bellman operator $\mathcal{T}$:
\begin{equation}
  (\mathcal{T}V)(a, z) = \max_{c, \ell} \left\{ u(c, \ell, G) + \beta_d \sum_{z'} \Pi(z, z') V(a', z') \right\}.
\end{equation}

Under our assumptions (bounded utility, $\beta_d < 1$, compact choice set), $\mathcal{T}$ is a contraction mapping with modulus $\beta_d$ in the sup-norm. By the Banach fixed-point theorem, $\mathcal{T}$ has a unique fixed point $V^*$ and the sequence $V^n = \mathcal{T}^n V^0$ converges to $V^*$ at rate $\beta_d^n$.

In practice, with $\beta_d = 0.95$, the error after $n$ iterations is bounded by $\beta_d^n / (1 - \beta_d) \cdot \|V^1 - V^0\|_\infty$. For our tolerance of $10^{-6}$, this typically requires 150--300 iterations depending on the state.

The discretization introduces approximation error from three sources:
\begin{enumerate}
  \item \textbf{Asset grid:} The exponential spacing provides higher resolution near $\underline{a}$ where marginal utility is highest. Linear interpolation between grid points introduces $O(\Delta a^2)$ error.
  \item \textbf{Leisure grid:} The uniform grid introduces at most $O(1/N_\ell)$ error in the leisure choice. The EGM solver uses 51 points (vs.\ 21 for VFI).
  \item \textbf{Consumption:} In VFI mode, an 11-point inner grid provides a coarse approximation. The EGM solver (\Cref{sec:egm}) eliminates this grid entirely by inverting the Euler equation, yielding exact consumption policy functions conditional on the asset grid.
\end{enumerate}


\section{Full Parameter Table}\label{app:params}

\Cref{tab:full_params} provides the complete parameter specification including derived and optional parameters.

\begin{table}[tp]
\centering
\caption{Complete Parameter Specification}\label{tab:full_params}
\small
\begin{tabular}{@{}llcl@{}}
\toprule
\textbf{Module} & \textbf{Parameter} & \textbf{Default} & \textbf{Constraint} \\
\midrule
Core & \texttt{num\_agents} & 50 & $\geq 20$ \\
Core & \texttt{max\_periods} & 100 & $\geq 1$ \\
Core & \texttt{seed} & 42 & integer \\
\midrule
Shocks & \texttt{rho\_z} & 0.9 & $(0, 1)$ \\
Shocks & \texttt{sigma\_z} & 0.2 & $> 0$ \\
Shocks & \texttt{num\_z\_states} & 7 & $[2, 50]$ \\
Shocks & \texttt{rho\_a} & 0.95 & $(0, 1)$ \\
Shocks & \texttt{sigma\_a} & 0.01 & $> 0$ \\
Shocks & \texttt{rho\_e} & 0.85 & $(0, 1)$ \\
Shocks & \texttt{sigma\_e} & 0.3 & $> 0$ \\
Shocks & \texttt{num\_e\_states} & 7 & $[2, 50]$ \\
Shocks & \texttt{enable\_preference\_shocks} & False & boolean \\
\midrule
Production & \texttt{alpha} & 0.33 & $(0, 1)$ \\
Production & \texttt{delta} & 0.1 & $(0, 1)$ \\
Production & \texttt{min\_firm\_capital} & 10.0 & $> 0$ \\
\midrule
Entrepreneurial & \texttt{firm\_entry\_cost} & 5.0 & $> 0$ \\
Entrepreneurial & \texttt{firm\_discount\_factor} & 0.95 & $(0, 1)$ \\
Entrepreneurial & \texttt{use\_bellman\_occ\_choice} & False & boolean \\
\midrule
Household & \texttt{a\_min} & 0.0 & $\geq 0$ \\
Household & \texttt{initial\_wealth\_mean} & 100.0 & $> 0$ \\
Household & \texttt{initial\_wealth\_std} & 30.0 & $\geq 0$ \\
Household & \texttt{homogeneous\_preferences} & True & boolean \\
Household & \texttt{utility\_alpha} & 0.4 & $(0, 1)$ \\
Household & \texttt{utility\_beta} & 0.35 & $(0, 1)$ \\
Household & \texttt{utility\_gamma} & 0.25 & $(0, 1)$ \\
Household & \texttt{utility\_beta\_discount} & 0.95 & $(0, 1)$ \\
\midrule
Market clearing & \texttt{market\_clearing\_method} & analytical & analytical/walrasian \\
Market clearing & \texttt{wage\_floor} & $10^{-6}$ & $> 0$ \\
Market clearing & \texttt{interest\_rate\_floor} & $-0.99$ & $> -1$ \\
\midrule
Intervals & \texttt{proposal\_interval} & 5 & $\geq 1$ \\
Intervals & \texttt{observer\_interval} & 5 & $\geq 1$ \\
\midrule
Governance & \texttt{default\_tax\_rate} & 0.1 & $[0, 1]$ \\
Governance & \texttt{proposal\_skip\_prob} & 0.85 & $[0, 1]$ \\
\midrule
LLM & \texttt{use\_llm} & True & boolean \\
LLM & \texttt{llm\_temperature} & 0.0 & $[0, 2]$ \\
LLM & \texttt{llm\_batch\_size} & 10 & $\geq 1$ \\
LLM & \texttt{llm\_cache\_enabled} & True & boolean \\
LLM & \texttt{benchmark\_mode} & False & boolean \\
\bottomrule
\end{tabular}
\end{table}

\begin{table}[tp]
\centering
\caption{Complete Parameter Specification (continued)}\label{tab:full_params2}
\small
\begin{tabular}{@{}llcl@{}}
\toprule
\textbf{Module} & \textbf{Parameter} & \textbf{Default} & \textbf{Constraint} \\
\midrule
R\&D & \texttt{rd\_success\_base\_prob} & 0.1 & $[0, 1]$ \\
R\&D & \texttt{rd\_tfp\_improvement\_mean} & 0.05 & $> 0$ \\
R\&D & \texttt{rd\_tfp\_improvement\_std} & 0.02 & $\geq 0$ \\
\midrule
Solver & \texttt{solver\_method} & egm & vfi/vfi\_numpy/egm \\
Distribution & \texttt{distribution\_mode} & individual & individual/kfe \\
Distribution & \texttt{kfe\_convergence\_tolerance} & $10^{-10}$ & $> 0$ \\
Distribution & \texttt{kfe\_max\_iterations} & 10000 & $\geq 1$ \\
\midrule
Two-asset & \texttt{two\_asset\_mode} & False & boolean \\
Two-asset & \texttt{chi\_0} & 0.01 & $\geq 0$ \\
Two-asset & \texttt{chi\_1} & 0.005 & $\geq 0$ \\
Two-asset & \texttt{b\_min} & 0.0 & --- \\
\midrule
Government & \texttt{initial\_debt} & 0.0 & $\geq 0$ \\
Government & \texttt{debt\_gdp\_max} & 1.5 & $> 0$ \\
Government & \texttt{fiscal\_rule\_adjustment} & 0.01 & $(0, 1)$ \\
\midrule
Nominal & \texttt{nominal\_rigidities} & False & boolean \\
Nominal & \texttt{rotemberg\_cost} & 100.0 & $\geq 0$ \\
Nominal & \texttt{taylor\_phi\_pi} & 1.5 & $> 1$ \\
Nominal & \texttt{taylor\_phi\_y} & 0.125 & $\geq 0$ \\
Nominal & \texttt{inflation\_target} & 0.02 & --- \\
Nominal & \texttt{elasticity\_sub} & 6.0 & $> 1$ \\
Nominal & \texttt{wage\_rigidity} & False & boolean \\
Nominal & \texttt{rotemberg\_wage\_cost} & 50.0 & $\geq 0$ \\
\midrule
Political & \texttt{political\_lambda} & 0.05 & $[0, 1]$ \\
Political & \texttt{pure\_bellman\_politics} & False & boolean \\
\bottomrule
\end{tabular}
\end{table}


\section{Sandbox Specification}\label{app:sandbox}

Constitutional rules may include enforcement code---Python expressions that are evaluated in a sandboxed environment to compute taxes, transfers, or other policy outcomes. The sandbox imposes strict safety constraints:

\textbf{Allowed AST nodes.} Only the following Python AST node types are permitted:
\begin{itemize}
  \item Literals (\texttt{Constant})
  \item Arithmetic operators (\texttt{Add}, \texttt{Sub}, \texttt{Mult}, \texttt{Div}, \texttt{FloorDiv}, \texttt{Mod}, \texttt{Pow})
  \item Unary operators (\texttt{USub}, \texttt{UAdd})
  \item Comparison operators (\texttt{Eq}, \texttt{NotEq}, \texttt{Lt}, \texttt{LtE}, \texttt{Gt}, \texttt{GtE})
  \item Boolean operators (\texttt{And}, \texttt{Or}, \texttt{Not})
  \item Variable references (\texttt{Name})
  \item Function calls (restricted to safe builtins)
  \item Conditional expressions (\texttt{IfExp})
  \item Simple assignments (\texttt{Assign})
  \item Container constructors (\texttt{Tuple}, \texttt{List}, \texttt{Dict})
\end{itemize}

\textbf{Disallowed constructs.} Imports, function/class definitions, attribute access, comprehensions, lambda expressions, \texttt{exec}, \texttt{eval}, and all forms of I/O are blocked at the AST level.

\textbf{Available builtins.} The sandbox provides a restricted set of built-in functions: \texttt{abs}, \texttt{min}, \texttt{max}, \texttt{round}, \texttt{sum}, \texttt{len}, plus mathematical functions \texttt{sqrt}, \texttt{log}, \texttt{exp}, and \texttt{pow}.

\textbf{Available variables.} Enforcement code for tax rules receives \texttt{income}, \texttt{rate}, and \texttt{wealth}. Transfer rules receive \texttt{revenue} and \texttt{num\_agents}. Public goods rules receive \texttt{fraction\_of\_revenue} and \texttt{revenue}.

This design ensures that constitutional rules, even when proposed by LLM agents, cannot compromise the simulation's integrity through arbitrary code execution.


% =============================================================================
% Bibliography
% =============================================================================
\bibliographystyle{plainnat}

\begin{thebibliography}{99}

\bibitem[Acemoglu(2005)]{acemoglu2005}
Acemoglu, D. (2005).
\newblock Politics and economics in weak and strong states.
\newblock \textit{Journal of Monetary Economics}, 52(7):1199--1226.

\bibitem[Acemoglu and Robinson(2006)]{acemoglurobinson2006}
Acemoglu, D. and Robinson, J.~A. (2006).
\newblock \textit{Economic Origins of Dictatorship and Democracy}.
\newblock Cambridge University Press.

\bibitem[Aiyagari(1994)]{aiyagari1994}
Aiyagari, S.~R. (1994).
\newblock Uninsured idiosyncratic risk and aggregate saving.
\newblock \textit{Quarterly Journal of Economics}, 109(3):659--684.

\bibitem[Bewley(1977)]{bewley1977}
Bewley, T. (1977).
\newblock The permanent income hypothesis: A theoretical formulation.
\newblock \textit{Journal of Economic Theory}, 16(2):252--292.

\bibitem[Cagetti and De~Nardi(2006)]{cagettideNardi2006}
Cagetti, M. and De~Nardi, M. (2006).
\newblock Entrepreneurship, frictions, and wealth.
\newblock \textit{Journal of Political Economy}, 114(5):835--870.

\bibitem[Huggett(1993)]{huggett1993}
Huggett, M. (1993).
\newblock The risk-free rate in heterogeneous-agent incomplete-insurance economies.
\newblock \textit{Journal of Economic Dynamics and Control}, 17(5):953--969.

\bibitem[Kopecky and Suen(2010)]{kopeckysuen2010}
Kopecky, K.~A. and Suen, R.~M. (2010).
\newblock Finite state {M}arkov-chain approximations to highly persistent processes.
\newblock \textit{Review of Economic Dynamics}, 13(3):701--714.

\bibitem[Krusell and Smith(1998)]{krusellsmith1998}
Krusell, P. and Smith, A.~A. (1998).
\newblock Income and wealth heterogeneity in the macroeconomy.
\newblock \textit{Journal of Political Economy}, 106(5):867--896.

\bibitem[Storesletten et~al.(2004)]{storesletten2004}
Storesletten, K., Telmer, C.~I., and Yaron, A. (2004).
\newblock Consumption and risk sharing over the life cycle.
\newblock \textit{Journal of Monetary Economics}, 51(3):609--633.

\bibitem[Carroll(2006)]{carroll2006}
Carroll, C.~D. (2006).
\newblock The method of endogenous gridpoints for solving dynamic stochastic optimization problems.
\newblock \textit{Economics Letters}, 91(3):312--320.

\bibitem[Fella(2014)]{fella2014}
Fella, G. (2014).
\newblock A generalized endogenous grid method for non-smooth and non-concave problems.
\newblock \textit{Review of Economic Dynamics}, 17(2):329--344.

\bibitem[Kaplan et~al.(2018)]{kaplanmollviolante2018}
Kaplan, G., Moll, B., and Violante, G.~L. (2018).
\newblock Monetary policy according to {HANK}.
\newblock \textit{American Economic Review}, 108(3):697--743.

\bibitem[Rotemberg(1982)]{rotemberg1982}
Rotemberg, J.~J. (1982).
\newblock Sticky prices in the {United States}.
\newblock \textit{Journal of Political Economy}, 90(6):1187--1211.

\bibitem[Taylor(1993)]{taylor1993}
Taylor, J.~B. (1993).
\newblock Discretion versus policy rules in practice.
\newblock \textit{Carnegie-Rochester Conference Series on Public Policy}, 39:195--214.

\bibitem[Young(2010)]{young2010}
Young, E.~R. (2010).
\newblock Solving the incomplete markets model with aggregate uncertainty using the {Krusell-Smith} algorithm and non-stochastic simulations.
\newblock \textit{Journal of Economic Dynamics and Control}, 34(1):36--41.

\end{thebibliography}

% =============================================================================
\end{document}
% =============================================================================
